\subsection{典范内射消解}

设 F 为 A-模 \(\mathrm{Hom}_{\mathbf{Z}}(\mathbf{A},\mathbf{Q} / \mathbf{Z})\) ；对于每个 A-模 M，我们设 \(\mathrm{I}^0 (\mathbf{M}) = \mathrm{F}^{\mathrm{Hom}_{\mathbf{A}}(\mathbf{M},\mathbf{F})}\) 并用 \(e_{\mathbf{M}}:\mathbf{M}\to \mathrm{I}^{0}(\mathbf{M})\) 表示将 \(m\in \mathbf{M}\) 映到族 \((\varphi (m))_{\varphi \in \mathrm{Hom}_{\mathbf{A}}(\mathbf{M},\mathbf{F})}\) 的同态。根据X，第19页，推论2， \(\mathrm{I}^0 (\mathbf{M})\) 是一个内射 A-模，且 \(e_{\mathbf{M}}\) 是内射的。设 \(\mathbf{K}^0 (\mathbf{M}) = \mathbf{Coker}\) \(e_{\mathbf{M}}\) 并记 \(q_{\mathbf{M}}:\mathrm{I}^0 (\mathbf{M})\to \mathrm{K}^0 (\mathbf{M})\) 为典范投影。因此我们有一个正合序列

\[
0 \longrightarrow \mathrm{M} \stackrel {e _ {\mathrm{M}}} {\longrightarrow} \mathrm{I} ^ {0} (\mathrm{M}) \stackrel {q _ {\mathrm{M}}} {\longrightarrow} \mathrm{K} ^ {0} (\mathrm{M}) \longrightarrow 0.
\]

我们定义一个分次A-模I(M)，设 \(\mathbf{I}^n (\mathbf{M}) = 0\) 于 \(n < 0\)，并对整数 \(n > 0\) 归纳地设

\[
\mathrm{I} ^ {n} (\mathbf {M}) = \mathrm{I} ^ {0} \left(\mathrm{K} ^ {n - 1} (\mathbf {M})\right), \quad \mathrm{K} ^ {n} (\mathbf {M}) = \mathrm{K} ^ {0} \left(\mathrm{K} ^ {n + 1} (\mathbf {M})\right).\tag{9}
\]

我们由下式定义 A-同态 \(\delta_{\mathbf{M}}^{n}:\mathrm{I}^{n}(\mathbf{M})\to \mathrm{I}^{n + 1}(\mathbf{M})\)

\[
\left\{ \begin{array}{l l} \delta_ {\mathbf {M}} ^ {n} = 0, & n <   0, \\ \delta_ {\mathbf {M}} ^ {0} = e _ {\mathbf {K} ^ {0} (\mathbf {M})} \circ q _ {\mathbf {M}}, \\ \delta_ {\mathbf {M}} ^ {n} = e _ {\mathbf {K} ^ {n} (\mathbf {M})} \circ q _ {\mathbf {K} ^ {n - 1} (\mathbf {M})}, & n > 0. \end{array} \right.\tag{10}
\]

由构造，我们有一个正合序列

\[
0 \longrightarrow \mathrm{M} \xrightarrow {e _ {\mathrm{M}}} \mathrm{I} ^ {0} (\mathrm{M}) \xrightarrow {\delta_ {\mathrm{M}} ^ {0}} \dots \longrightarrow \mathrm{I} ^ {n} (\mathrm{M}) \xrightarrow {\delta_ {\mathrm{M}} ^ {n}} \mathrm{I} ^ {n + 1} (\mathrm{M}) \longrightarrow \dots ,
\]

因此，如果把 \(e_{M}\) 延拓为复形的态射

\[
e _ {\mathbf {M}}: \mathbf {M} \rightarrow (\mathrm{I} (\mathbf {M}), \delta_ {\mathbf {M}}),
\]

我们得到M的一个内射分解，称为M的典范内射分解。设 \(f: \mathbf{M} \to \mathbf{N}\) 是 A-模的同态。用 \(\mathrm{I}^{0}(f)\) 表示从 \(\mathrm{I}^{0}(\mathbf{M}) = \mathrm{F}^{\mathrm{Hom}_{\mathbf{A}}(\mathbf{M},\mathbf{F})}\) 到 \(\mathrm{I}^{0}(\mathbf{N}) = \mathrm{F}^{\mathrm{Hom}_{\mathbf{A}}(\mathbf{N},\mathbf{F})}\) 的、将族 \((x_{\varphi})_{\varphi \in \mathrm{Hom}_{\mathbf{A}}(\mathbf{M},\mathbf{F})}\) 映到族 \((x_{\psi \circ f})_{\psi \in \mathrm{Hom}_{\mathbf{A}}(\mathbf{N},\mathbf{F})}\) 的同态。我们有：

\[
\mathrm{I} ^ {0} (f) \circ e _ {\mathbf {M}} = e _ {\mathbf {N}} \circ f.\tag{11}
\]

因此， \(\mathrm{I}^0 (f)\) 诱导一个同态 \(\mathbf{K}^0 (f):\mathbf{K}^0 (\mathbf{M})\to \mathbf{K}_-^0 (\mathbf{N})\) 并且我们有

\[
\mathrm{K} ^ {0} (f) \circ q _ {\mathbf {M}} = q _ {\mathbf {N}} \circ \mathrm{K} ^ {0} (f).\tag{12}
\]

令 \(\mathrm{I}^n(f) = 0\) 当 \(n < 0\) 并通过对整数 \(n > 0\) 的归纳，定义同态 \(\mathrm{I}^n(f): \mathrm{I}^n(\mathbf{M}) \to \mathrm{I}^n(\mathbf{N})\) 与 \(\mathrm{K}^n(f): \mathrm{K}^n(\mathbf{M}) \to \mathrm{K}^n(\mathbf{N})\) 如下：

\[
\left\{ \begin{array}{l} \mathrm{I} ^ {n} (f) = \mathrm{I} ^ {0} (\mathrm{K} ^ {n - 1} (f)) \\ \mathrm{K} ^ {n} (f) = \mathrm{K} ^ {0} (\mathrm{K} ^ {n - 1} (f)). \end{array} \right.\tag{13}
\]

命题5. --- I(f) : I(M) → I(N) 是A-模复形的一个态射；我们有 I(f) ∘ e\_M = e\_N ∘ f。

这是以与命题4类似的方式证明的。

我们有

\[
\mathrm{I} (1 _ {\mathbf {M}}) = 1 _ {\mathrm{I} (\mathbf {M})}\tag{14}
\]

且对于每个A-模同态 \(g: N \rightarrow P\)

\[
\mathrm{I} (g \circ f) = \mathrm{I} (g) \circ \mathrm{I} (f).\tag{15}
\]

注记. --- 如果 \(f, g \in \operatorname{Hom}_{\mathbf{A}}(\mathbf{M}, \mathbf{N})\) ，我们没有 \(\mathrm{I}(f + g) = \mathrm{I}(f) + \mathrm{I}(g)\)。然而，这两个态射依X，第49页，命题3 bis是同伦的。

若M是一个右A-模，我们设 \(\mathrm{I}(\mathbf{M}) = \mathrm{I}(\mathbf{M}^{\circ})^{\circ}\) ；它被称为M的典范内射分解，并且我们有

\[
\mathrm{I} (\mathbf {M} ^ {\circ}) = \mathrm{I} (\mathbf {M}) ^ {\circ}.
\]

\subsection{有限生成分解}

特别地，从前两节可以得出，每个A-模都有内射分解、自由分解（因此也有投射分解或平坦分解）。在某些情况下，可以更精确地表述：

假设A是左诺特环，令M为A-模。让我们递归地构造序列 \((\mathbf{L}_{n})_{n\geqslant0},(\mathbf{Z}_{n})_{n\geqslant0},(d_{n})_{n\geqslant1}\) 其中，对所有 \(n\geqslant0\)，\(L_{n}\) 是有限生成的自由A-模，\(Z_{n}\) 是 \(L_{n}\) 的一个子模，而 \(d_{n+1}:L_{n+1}\to L_{n}\) 是一个同态。为此，选取 M 的一个有限生成族 \((m_{i})_{i\in I_{0}}\)，设 \(L_{0}=A^{(I_{0})}\)，定义 \(p:L_{0}\to M\) 为 \(p(e_{i})=m_{i}\)，并设 \(Z_{0}=\operatorname{Ker}(p)\)。对于 \(n\geqslant0\)，设模 \(L_{n}\) 与 \(Z_{n}\) 已经构造完成；由于 \(Z_{n}\) 含于 \(L_{n}\)，它是有限型的；选取 \((x_{n,i})_{i\in I_{n+1}}\) 作为 \(Z_{n}\) 的一个有限生成族；令 \(L_{n+1}=A^{(I_{n+1})}\)，定义 \(d_{n+1}\) 为 \(d_{n+1}(e_{i})=x_{n,i}\)，并设 \(Z_{n+1}=\operatorname{Ker}(d_{n+1})\)。根据构造，我们有一个正合序列

\[
\dots \longrightarrow \mathrm{L} _ {n + 1} \xrightarrow {d _ {n + 1}} \mathrm{L} _ {n} \longrightarrow \dots \longrightarrow \mathrm{L} _ {0} \xrightarrow {p} \mathrm{M} \longrightarrow 0,
\]

由此得出：

命题6. --- 当A是左诺特环时，每个有限生成的A-模M都有一个自由分解 \(p: \mathbf{L} \to \mathbf{M}\)，使得 \(\mathbf{L}_n\) 对每个整数 \(n\) 都是有限生成的。更一般地：

命题7. --- 设C是一个A-复形，且设 \(a \in \mathbf{Z}\) 使得 \(\mathrm{H}_{n}(\mathbf{C}) = 0\) 对 \(n < a\) 成立。a) 存在一个自由A-复形 \(\mathbf{L}\)，使得 \(\mathbf{L}_{n} = 0\) 对 \(n < a\) 成立，以及一个拟同构 \(f: \mathbf{L} \to \mathbf{C}\)。

\begin{enumerate}
\def\labelenumi{\alph{enumi})}
\setcounter{enumi}{1}
\tightlist
\item
  假设A是左诺特的，且A-模 \(\mathrm{H}_{n}(\mathbf{C})\) , \(n \in Z\) 是有限生成的。存在一个自由A-复形L，使得 \(L_{n} = 0\) 对 n \textless{} a 成立，且 \(L_{n}\) 对每个 n 都是有限生成的 A-模，并且存在一个拟同构 \(f: L \to C\)。
\end{enumerate}

设 C' 为 C 的子复形，使得 \(C_{n}^{\prime}=C_{n}\) 对于 n\textgreater a, \(\mathbf{C}_{a}^{\prime}=\mathbf{Z}_{a}(\mathbf{C})\) , \(C_{n}^{\prime}=0\) 对于 n \textless{} a；则从 C' 到 C 的典范单射是一个拟同构。将 C 替换为 \(C'\)，我们因此可以假设 \(C_{n}=0\) 对于 n\textless a。于是该命题可由以下引理对 r=a，\(a+1,\ldots\) 的反复应用得出：

引理3. --- 设 C 是一个复形且 \(r \in \mathbf{Z}\)。存在一个复形 \(\mathbf{C}'\) 以及一个拟同构 \(f: \mathbf{C}' \to \mathbf{C}\)，使得 \(f_n: \mathbf{C}_n' \to \mathbf{C}_n\) 对 \(n < r\) 是同构，并且 \(\mathbf{C}_r'\) 是自由A-模。若A是诺特环且A-模 \(\mathrm{H}_r(\mathbf{C})\) 与 \(\mathbf{C}_{r-1}\) 是有限生成的，则可要求 \(\mathbf{C}_r'\) 是有限生成的。

\begin{enumerate}
\def\labelenumi{\alph{enumi})}
\item
  首先设 \(h: \mathbf{M} \to \mathbf{C}_r\) 为A-模的同态；记 \(d = (d_n)\) 为C的微分。设N为 \(\mathbf{M} \times \mathbf{C}_{r+1}\) 中由对 \((m, x)\)（满足 \(h(m) = d_{r+1}(x)\)）组成的子模；我们定义复形 \((\mathbf{C}', d')\) 如下：\(\mathbf{C}_n' = \mathbf{C}_n\) 对 \(n \neq r, r+1, \mathbf{C}_r' = \mathbf{M}, \mathbf{C}_{r+1}' = \mathbf{N}, d_n' = d_n\)，对 \(n \neq r, r+1, r+2, d_r' = d_r \circ h, d_{r+1}'(m, x) = m\) 成立，其中 \((m, x) \in \mathbf{N}\)，而 \(d_{r+2}'(y) = (0, d_{r+2}(y))\) 对 \(y \in \mathbf{C}_{r+2}\) 成立。再考虑复形的态射 \(f: \mathbf{C}' \to \mathbf{C}\)：\(f_n = 1_{\mathbf{C}_n}\) 对 \(n \neq r, r+1, f_r = h, f_{r+1}(m, x) = x\) 成立。
\item
  复形Ker f在次数 \(\neq r, r + 1\) 处为零，且微分 \(d_{r+1}^{\prime}\) 诱导出从 \(\operatorname{Ker} f_{r+1}\) 到 \(\operatorname{Ker} f_r\) 上的同构，因此 \(\operatorname{H}(\operatorname{Ker} f) = 0\)。
\item
  当复合映射 \(\mathbf{M} \xrightarrow{\hbar} \mathbf{C}_r \to \mathbf{C}_r / \mathbf{B}_r(\mathbf{C})\) 是满射时，类似地可看出 \(\mathrm{H}(\mathrm{Coker}f) = 0\) ，于是 \(f\) 是一个拟同构（X，第31页，推论2）。
\item
  当假设A是诺特环且A-模 \(H_{r}(C)\) 与 \(C_{r-1}\) 是有限生成的，则 \(C_{r}/B_{r}(C)\) 是有限生成的，依据正合序列（X，第25页）
\end{enumerate}

\[
0 \to \mathrm{H} _ {r} (\mathrm{C}) \to \mathrm{C} _ {r} / \mathrm{B} _ {r} (\mathrm{C}) \to \mathrm{C} _ {r - 1};
\]

于是存在有限生成的自由A-模M及一个同态 \(h: \mathbf{M} \to \mathbf{C}_{r}\) 使得 \(c)\) 的条件成立；在一般情况下，存在自由模M及一个满同态 \(h: \mathbf{M} \to \mathbf{C}_{r}\) 。这就完成了证明。

\subsection{极小投射分解}

设 M 是一个 A-模，且设

\[
\dots \longrightarrow \mathrm{P} _ {n} \xrightarrow {d _ {n}} \mathrm{P} _ {n - 1} \longrightarrow \dots \longrightarrow \mathrm{P} _ {0} \xrightarrow {d _ {0}} \mathrm{M} \longrightarrow 0\tag{P}
\]

为M的一个分解。我们称(P)为极小投射分解，如果对每一个 \(n \geqslant 0\)，同态 \(\delta_{n}: P_{n} \to \operatorname{Im}(d_{n})\)（由 \(d_{n}\) 所诱导）是一个投射覆盖（VIII，§8，第5号）。

命题8.—设M是一个A-模，且P与P'是M的两个极小投射分解，其中f：P→P'是分解的一个态射。则f是一个同构。特别地，M的任意两个极小投射分解都是同构的。

令 \(\tilde{\mathbf{P}}_{n} = \mathbf{P}_{n}\) 于 \(n \neq -1\)，且 \(\tilde{\mathbf{P}}_{-1} = \mathbf{M}\)；让我们类似地定义 \(\tilde{\mathbf{P}}_{n}^{\prime}\) 并设 \(f_{-1} = 1_{\mathbf{M}}\)。让我们从 \(-1\) 出发用归纳法证明 \(f_{n}: \tilde{\mathbf{P}}_{n} \to \tilde{\mathbf{P}}_{n}^{\prime}\) 对所有 \(n\) 都是同构。这对 \(n = -1\) 是显然的；假设 \(f_{n}\) 与 \(f_{n-1}\) 是同构。由图的交换性可得：

\[
\begin{array}{c c c} \mathbf {P} _ {n} & \xrightarrow {d _ {n}} & \mathbf {P} _ {n - 1} \\ \Big \downarrow f _ {n} & & \Big \downarrow f _ {n - 1} \\ \mathbf {P} _ {n} ^ {\prime} & \xrightarrow {d _ {n} ^ {\prime}} & \mathbf {P} _ {n - 1} ^ {\prime} \end{array}
\]

即 \(f_{n}\) 诱导出 \(g_{n}\)，它是从 \(\operatorname{Ker} d_{n}\) 到 \(\operatorname{Ker} d_{n}^{\prime}\) 上的一个同构。于是由图的交换性以及

\[
\begin{array}{c c c} \mathrm{P} _ {n + 1} & \xrightarrow {\delta_ {n + 1}} & \text {Ker} d _ {n} \\ \Big \downarrow f _ {n + 1} & & \Big \downarrow g _ {n} \\ \mathrm{P} _ {n + 1} ^ {\prime} & \xrightarrow {\delta_ {n + 1} ^ {\prime}} & \text {Ker} d _ {n} ^ {\prime} \end{array}
\]

以及 VIII，前引文献，可得 \(f_{n+1}\) 是一个同构。

推论. --- 设 M 是一个 A-模，P 和 P' 是 M 的两个投射分解；假设 P 是极小的。设 f : P → P' 和 g : P' → P 是两个分解态射。则 f 是单射，g 是满射，且 P' 是子复形 Im f 与 Ker g 的直和。此外，Ker g 具有零同调。

事实上， \(\alpha = g \circ f\) 是P的一个自同构（命题8）。设 \(\tilde{f} = f \circ \alpha^{-1}\) 。我们有

\[
\operatorname{Im} \tilde {f} = \operatorname{Im} f \quad \text { and } \quad g \circ \tilde {f} = 1 _ {\mathrm{P}},
\]

这表明 \(\mathbf{P}' = \operatorname{Im} \tilde{f} \oplus \operatorname{Ker} g\) 。由于序列

\[
0 \rightarrow \operatorname{Ker} g \rightarrow \mathrm{P} ^ {\prime} \stackrel {g} {\rightarrow} \mathrm{P} \rightarrow 0
\]

是正合的，且 \(g\) 是一个拟同构， \(\operatorname{Ker} g\) 具有零同调。

命题9。——设M是一个A-模，且(P, p)是M的一个投射分解。设r表示A的根。假设以下任一条件成立： \(P_{n}\) 对所有n都是有限生成的A-模，或者r是幂零的。那么 (P, p) 为极小的充要条件是复形 \((\mathbf{A}/\mathbf{r}) \otimes_{\mathbf{A}} \mathbf{P}\) 具有零微分，换言之即

\[
d _ {n + 1} (\mathbf {P} _ {n + 1}) \subset \mathrm{r} \mathbf {P} _ {n} \quad \text { for all } n \geqslant 0.
\]

假设 (P, p) 是极小的。根据 VIII，前引文献，同态

\[
1 \otimes \delta_ {n}: (\mathrm{A} / \mathrm{r}) \otimes_ {\mathrm{A}} \mathrm{P} _ {n} \rightarrow (\mathrm{A} / \mathrm{r}) \otimes_ {\mathrm{A}} \operatorname{Im} d _ {n}
\]

是一个同构。于是由正合序列可得

\[
0 \longrightarrow \operatorname{Im} d _ {n + 1} \xrightarrow {j _ {n}} \mathrm{P} _ {n} \xrightarrow {\delta_ {n}} \operatorname{Im} d _ {n} \longrightarrow 0
\]

使得同态 \(1 \otimes j_n : (\mathbf{A} / \mathfrak{r}) \otimes_{\mathbf{A}} \operatorname{Im} d_{n+1} \to (\mathbf{A} / \mathfrak{r}) \otimes_{\mathbf{A}} \mathrm{P}_n\) 为零；由于 \(d_{n+1} = j_n \circ \delta_{n+1}\) ，我们推出 \(1_{\mathbf{A} / \mathfrak{r}} \otimes d_{n+1} = 0\) ，对 \(n \geqslant 0\) 。

反之，假设对所有 \(n \geqslant 1\) , \(1 \otimes d_n\) 为零，即 \(\operatorname{Im} d_n = \operatorname{Ker} d_{n-1}\) 包含于 \(\mathbf{rP}_{n-1}\) 。由于 \(\delta_{n-1}\) 是满射的，由VIII，前引文献可知 \(\delta_{n-1}\) 对 \(n \geqslant 1\) 是一个投射覆盖，因此 \((\mathbf{P}, p)\) 是极小的。

命题10。——设A是左诺特局部环，令M为有限生成的A模。则M具有极小分解(P, p)；对每个 \(n \geqslant 0\) , \(\mathbf{P}_{n}\) 是有限生成的自由模。

事实上，在第5节（第53页）所作的构造中，由VIII，前引文献，可以取 \((\mathrm{L}_{0}, p)\) 为 M 的一个投射覆盖，并取 \(L_{n+1}\) 为 Ker \(d_{n}\) 的一个投射覆盖。如此得到的分解便是极小的。

注记. --- 1) 设 m 表示 A 的极大理想，并设 \(k = \mathbf{A}/\mathfrak{m}\) 。设 P 为 M 的一个极小投射分解，并设 \(b_{n} = \dim_{k}(k \otimes_{\mathbf{A}} \mathbf{P}_{n})\) 。则 \(\mathbf{P}_{n}\) 是秩为 \(b_{n}\) 的自由 A-模。由命题 8 的推论可知，对于 M 的任何其他投射分解 \(\mathbf{P}'\) ，有 \(\dim_{k}(k \otimes_{\mathbf{A}} \mathbf{P}_{n}') \geqslant b_{n}\) ，且等号成立当且仅当 \(\mathbf{P}'\) 是极小的。

\begin{enumerate}
\def\labelenumi{\arabic{enumi})}
\setcounter{enumi}{1}
\tightlist
\item
  由命题 9， \(b_n\) 是 \(k\) 上 \(H_n(k \otimes_A P)\) 的维数，* 换言之即 \(\text{Tor}_n^A(k, M)\) 的维数。它也是 \(k\) 上 \(\text{Ext}_A^n(M, k)\) 的维数（参见 X，第 103 页，注记 3）*。
\end{enumerate}

\subsection{分次分解}

在本节中，假设环 A 配备了一个分次结构 \((\mathbf{A}_n)_{n\in \mathbf{Z}}\) ，使得 \(\mathbf{A}_n = 0\) 当 \(n < 0\) 。称一个分次 A-模 M 是下有界的，如果 \(\mathbf{M}_n = 0\) 当n足够小时；每个有限生成的分次A-模都是下有界的。

命题11. --- 若 M 是下方有界的分次 A-模（相应地，若 M 是有限生成的分次 A-模且 A 是左诺特的），则存在一个分次 A-模的正合序列，向左无界，

\[
\dots \longrightarrow \mathrm{L} _ {n} \xrightarrow {d _ {n}} \mathrm{L} _ {n - 1} \longrightarrow \dots \longrightarrow \mathrm{L} _ {1} \xrightarrow {d _ {1}} \mathrm{L} _ {0} \xrightarrow {d _ {0}} \mathrm{M} \longrightarrow 0
\]

其中 \(L_{i}\) 是分次自由的且有下界（相应地，分次自由且有限生成），并且其中 \(d_{i}\) 是次数为 0 的分次同态。

若 N 是下方有界的分次 A-模（相应地，且在诺特环 A 上有限生成），则存在下方有界（相应地，且有限生成）的分次自由 A-模 L 以及一个满射的分次同态 L → N（II，第 167 页，注记 3）。

既然如此，假设给定一个分次A-模与次数为0的分次同态的正合序列

\[
\mathrm{L} _ {n} \xrightarrow {d _ {n}} \mathrm{L} _ {n - 1} \longrightarrow \dots \longrightarrow \mathrm{L} _ {0} \xrightarrow {d _ {0}} \mathrm{M} \longrightarrow 0,
\]

其中 \(L_{i}, i = 0, \ldots, n\) ，它们是分次自由的且有下界（分别为分次自由且有限生成的）。则N = Ker \(d_{n}\) 下有界（相应地，有限生成）；因此存在一个有下界的分次自由A-模（相应地，分次自由且有限生成） \(L_{n+1}\) 以及一个次数为0的分次同态 \(d_{n+1}: L_{n+1} \to L_{n}\)，使得Im \(d_{n+1} = N\)；序列

\[
\mathrm{L} _ {n + 1} \xrightarrow {d _ {n + 1}} \mathrm{L} _ {n} \xrightarrow {d _ {n}} \mathrm{L} _ {n - 1} \longrightarrow \dots \longrightarrow \mathrm{L} _ {0} \xrightarrow {d _ {0}} \mathrm{M} \longrightarrow 0
\]

于是正合。该命题随后由前述内容对n归纳得出。

\subsection{标准分解}

在本节中，假设环A是交换环 \(k\) 上的（结合且有单位的）代数。对于 \(n \geqslant 0\)，我们用 \(\mathbf{B}_n\) 表示在 \(k\) 上的 \((n + 2)\) 个都等于 A 的模的张量积。我们将其视为 (A, A)-双模，赋予其左（相应地，右）A-模结构，该结构由张量积的第一个（相应地，最后一个）因子的左（相应地，右）A-模结构诱导。

对于 \(n \geqslant 1\)，我们定义双模同态 \(d_n^i\)（对于 \(0 \leqslant i \leqslant n\)）和 \(d_n\)，它们从 \(B_n\) 到 \(B_{n-1}\)，由公式：

\[
d _ {n} ^ {i} (x _ {0} \otimes \ldots \otimes x _ {n + 1}) = x _ {0} \otimes \ldots \otimes x _ {i} x _ {i + 1} \otimes \ldots \otimes x _ {n + 1}, \quad 0 \leqslant i \leqslant n,
\]

\[
d _ {n} = \sum_ {i = 0} ^ {n} (- 1) ^ {i} d _ {n} ^ {i}.
\]

显然，

\[
d _ {n - 1} ^ {i} \circ d _ {n} ^ {j} = d _ {n - 1} ^ {j - 1} \circ d _ {n} ^ {i} \quad \text { for } \quad i <   j
\]

并且因此

\[
d _ {n - 1} \circ d _ {n} = \sum_ {0 \leqslant i <   j \leqslant n} (- 1) ^ {i + j} d _ {n - 1} ^ {i} \circ d _ {n} ^ {j} + \sum_ {0 \leqslant j \leqslant i \leqslant n - 1} (- 1) ^ {i + j} d _ {n - 1} ^ {i} \circ d _ {n} ^ {j} = 0.
\]

因此，如果我们设 \(\mathbf{B}_n = 0\) 于 \(n < 0\)，且 \(d_n = 0\) 于 \(n \leqslant 0\)，则序列 \((\mathbf{B}_n, d_n)\) 定义了 (A, A)-双模的复形（X，第 43 页），将其记为 \(\mathbf{B}(\mathbf{A})\)。对于每个左 A-模 \(\mathbf{M}\)，我们用 \(\mathbf{B}(\mathbf{A}, \mathbf{M})\) 表示由 \(\mathbf{B}_n \otimes_{\mathbf{A}} \mathbf{M}\) 和 \(d_n \otimes 1_{\mathbf{M}}\) 构成的复形，换言之即张量积复形 \(\mathbf{B}(\mathbf{A}) \otimes_{\mathbf{A}} \mathbf{M}_{*}\)；它是一个左A-模复形。

我们定义一个A-线性映射 \(\varepsilon_{\mathbf{M}}\)，它从 \(\mathbf{B}_{0}(\mathbf{A},\mathbf{M}) = (\mathbf{A}\otimes_{k}\mathbf{A})\otimes_{\mathbf{A}}\mathbf{M}\) 到 \(\mathbf{M}\)，由公式 \(\varepsilon_{\mathbf{M}}(a\otimes b\otimes m) = abm\)（\(a,b\in \mathbf{A},m\in \mathbf{M}\)）给出。我们有 \(\varepsilon_{\mathbf{M}}\circ d_1 = 0\) ，使得分次同态 \(\overline{\varepsilon}_{\mathbf{M}}:\mathbf{B}(\mathbf{A},\mathbf{M})\to \mathbf{M}\) ，它在次数 0 上与 \(\varepsilon_{\mathbf{M}}\) 一致，是 A-模复形的一个态射。

命题 12. —— 映射 \(\overline{\varepsilon}_{\mathbf{M}}: \mathbf{B}(\mathbf{A}, \mathbf{M}) \to \mathbf{M}\) 是 \(k\) -模复形的一个同伦等价。特别地，复形 \(\mathbf{B}(\mathbf{A}, \mathbf{M})\) 在 \(k\) 上是分裂的，且 \((\mathbf{B}(\mathbf{A}, \mathbf{M}), \overline{\varepsilon}_{\mathbf{M}})\) 是A-模M的一个左分解。

对于 \(n \geqslant 0\) ，让我们定义一个 \(k\) -线性映射 \(s_n: \mathbf{B}_n \to \mathbf{B}_{n+1}\) 通过以下公式：

\[
s _ {n} \left(x _ {0} \otimes \dots \otimes x _ {n + 1}\right) = 1 \otimes x _ {0} \otimes \dots \otimes x _ {n + 1} \quad \text { for } \quad x _ {0}, \dots , x _ {n + 1} \in A.
\]

它是右A-模的同态，满足以下恒等式：

\[
d _ {n + 1} ^ {i} \circ s _ {n} = s _ {n - 1} \circ d _ {n} ^ {i - 1} \quad \mathrm{for} \quad n \geqslant 1, \quad 1 \leqslant i \leqslant n + 1,
\]

\[
d _ {n + 1} ^ {0} \circ s _ {n} = 1 _ {\mathrm{B} _ {n}} \quad \text { for } \quad n \geqslant 1,
\]

因此我们有

\[
d _ {n + 1} \circ s _ {n} + s _ {n - 1} \circ d _ {n} = 1 _ {\mathrm{B} _ {n}} \quad \text { for } \quad n \geqslant 1.\tag{16}
\]

此外，我们有

\[
d _ {1} \circ s _ {0} (x _ {0} \otimes x _ {1}) = x _ {0} \otimes x _ {1} - 1 \otimes x _ {0} x _ {1} \quad \text { for } \quad x _ {0}, x _ {1} \in A.\tag{17}
\]

我们用 \(\eta: \mathbf{A} \to \mathbf{A} \otimes_k \mathbf{A}\) 表示由 \(\eta(a) = 1 \otimes a\) 定义的映射，且用 \(\overline{\eta}: \mathbf{A} \to \mathbf{B}(\mathbf{A})\) 表示与 \(\eta\) 在次数0上一致的复形的态射。显然 \(\overline{\varepsilon}_{\mathbf{A}} \circ \overline{\eta} = 1_{\mathbf{A}}\) ；公式（16）和（17）表明 \(d \circ s + s \circ d = 1_{\mathrm{B(A)}} - \overline{\eta} \circ \overline{\varepsilon}_{\mathbf{A}}\) 。设 \(\overline{\eta}_{\mathrm{M}} = \overline{\eta} \otimes 1_{\mathrm{M}}\) , \(d_{\mathrm{M}} = d \otimes 1_{\mathrm{M}}\) 与 \(s_{\mathrm{M}} = s \otimes 1_{\mathrm{M}}\)，我们推出 \(\overline{\varepsilon}_{\mathrm{M}} \circ \overline{\eta}_{\mathrm{M}} = 1_{\mathrm{M}}\) 与 \(d_{\mathrm{M}} \circ s_{\mathrm{M}} + s_{\mathrm{M}} \circ d_{\mathrm{M}} = 1_{\mathrm{B(A,M)}} - \overline{\eta}_{\mathrm{M}} \circ \overline{\varepsilon}_{\mathrm{M}}\)。换言之（X，第33页，定义5）， \(\overline{\varepsilon}_{\mathrm{M}}\) 是 \(k\) -模复形的一个同伦等价。命题的其他断言立即成立。

定义3. --- 左分解(B(A, M), \(\overline{\varepsilon}_{\mathrm{M}}\) ) 称为A-模M的标准分解。

若A和M是投射（分别为自由、平坦）k-模，则标准分解B(A, M)是M的投射（分别为自由、平坦）分解。

\subsection{分解与Grothendieck群}

若 \(\mathcal{C}\) 是由A-模的类组成的一个集合，我们将称左分解 \((\mathbf{P}, p)\) 是 \(\mathcal{C}\) 型有界的，如果复形 \(\mathbf{P}\) 是 \(\mathcal{C}\) 型有界的 (X，第41页)。

定理1. --- 设 \(\mathcal{C}_{0}\) 与 \(\mathcal{C}\) 是两个加性且左正合的A-模类，使得 \(\mathcal{C}_{0} \subset \mathcal{C}\)，并且每个 \(\mathcal{C}\) 型的A-模都具有 \(\mathcal{C}_{0}\) 型的有界左分解。则同态 \(\alpha : \mathrm{K}(\mathcal{C}_{0}) \to \mathrm{K}(\mathcal{C})\)（由 \(\mathcal{C}_{0}\) 到 \(\mathcal{C}\) 的包含关系导出的）是双射；若 M 是一个 \(\mathcal{C}\) 型的A-模，且 P 是 M 的一个 \(\mathcal{C}_{0}\) 型有界左分解，则我们有 \(\alpha^{-1}([M]_{\mathcal{C}}) = \chi_{\mathcal{C}_{0}}(P)\) （X，第 41 页，例 6）。

引理4. --- 设 \(f: \mathbf{M}' \to \mathbf{M}\) 是 \(\mathcal{C}\) 型A-模的同态，且 \(p: \mathbf{P} \to \mathbf{M}\) 是 \(\mathbf{P}\) 的一个 \(\mathcal{C}_0\) 型有界的左分解。存在一个左分解 \(p': \mathbf{P}' \to \mathbf{M}'\)（\(\mathcal{C}_0\) 型有界的）以及一个复形的态射 \(u: \mathbf{P}' \to \mathbf{P}\) ，使得 \(p \circ u = f \circ p'\)。

让我们对P的长度n进行归纳推理，当该长度为 \textless{} 0时命题是平凡的。考虑映射 \(g: \mathbf{M}' \times \mathbf{P}_0 \to \mathbf{M}\) 使得

\[
g (x, y) = f (x) - p _ {0} (y) \quad \text { for } \quad x \in \mathbf {M} ^ {\prime}, y \in \mathbf {P} _ {0},
\]

及其核K；由于g是满射且 \(M' \times P_{0}\) 和M是C类型的，A-模K是C类型的。设 \(h : P_{0}' \to K\) 是一个满射同态，其中 \(P_{0}'\) 是 \(C_{0}\) 型的；记 \(p_{0}' : P_{0}' \to M'\)（相应地， \(u_{0} : P_{0}' \to P_{0}\)）为 h 与投影 \(K \to M\)（相应地， \(K \to P_{0}\)）的复合同态；同态 \(p_{0}'\) 是满射的，并且我们有一个交换图

\[
\begin{array}{c} \mathbf {P} _ {0} ^ {\prime} \xrightarrow {u _ {0}} \mathbf {P} _ {0} \\ p _ {0} ^ {\prime} \Big \downarrow \qquad \qquad \qquad \qquad \qquad \qquad \qquad \qquad \qquad \Big \downarrow p _ {0} \\ \mathbf {M} ^ {\prime} \xrightarrow [ f ]{} \mathbf {M}. \end{array}
\]

然后只需将归纳假设应用于同态

\[
\operatorname{Ker} p _ {0} ^ {\prime} \rightarrow \operatorname{Ker} p _ {0}
\]

由 \(u_{0}\) .

引理5. --- 考虑一个交换图

\[
\begin{array}{c} \mathbf {P} ^ {\prime} \xrightarrow {u} \mathbf {P} \\ p ^ {\prime} \Big \downarrow \quad \Big \downarrow p \\ 0 \longrightarrow \mathbf {M} ^ {\prime} \xrightarrow {f} \mathbf {M} \longrightarrow \mathbf {M} ^ {\prime \prime} \longrightarrow 0 \end{array}
\]

其中 (P, p)（相应地 (P′, p′)）是 M（相应地 M′）的左分解，且底部水平行是一个正合序列。存在一个拟同构 \(p''\) : Con \((u) \to \mathbf{M}''\) 。

事实上，正合序列 (X，第 37 页，命题 7)

\[
0 \to \mathrm{P} \xrightarrow {\pi} \operatorname{Con} (u) \xrightarrow {\delta} \mathrm{P} ^ {\prime} (- 1) \to 0
\]

给出一个正合同调序列

\[
\begin{array}{c}\rightarrow \mathrm{H} _ {n} (\mathrm{P}) \rightarrow \mathrm{H} _ {n} (\mathrm{Con} (u)) \rightarrow \mathrm{H} _ {n - 1} (\mathrm{P} ^ {\prime}) \rightarrow \dots\\\dots \rightarrow \mathrm{H} _ {1} (\mathrm{Con} (u)) \rightarrow \mathrm{H} _ {0} (\mathrm{P} ^ {\prime}) \stackrel {{\partial}} {{\rightarrow}} \mathrm{H} _ {0} (\mathrm{P}) \rightarrow \mathrm{H} _ {0} (\mathrm{Con} (u)) \rightarrow 0.\end{array}
\]

根据 X，第 38 页，引理 3 a)，我们有 \(\partial = -\mathrm{H}_{0}(u)\) 。由于 \(\mathrm{H}_{n}(\mathbf{P}) = 0 = \mathrm{H}_{n}(\mathbf{P}')\)（\(n > 0\)），且 \(\mathrm{H}_{0}(u): \mathrm{H}_{0}(\mathrm{P}') \to \mathrm{H}_{0}(\mathbf{P})\) 与 \(f: \mathbf{M}' \to \mathbf{M}\) 等同，我们得出结论 \(\mathrm{H}_{n}(\mathrm{Con}(u)) = 0\)（对 \(n > 0\)），并且 \(\mathrm{H}_{0}(\mathrm{Con}(u))\) 同构于 \(\mathbf{M}''\) ，从而引理得证。

现在让我们证明该定理。

\begin{enumerate}
\def\labelenumi{\alph{enumi})}
\tightlist
\item
  设 M 是一个 C 型的 A-模。对于 M 的每一个左分解 \((\mathbf{P}, p)\)（\(C_{0}\) 型有界的），\(\chi_{\mathcal{C}_{0}}(\mathbf{P})\) 作为 \(\mathbf{K}(\mathcal{C}_{0})\) 中的元素仅依赖于 M。事实上，设 \((\mathbf{P}_{1}, p_{1})\) 与 \((\mathbf{P}_{2}, p_{2})\) 是该类型的两个分解。考虑 A-模
\end{enumerate}

\[
(\mathbf {P} _ {1} \times \mathbf {P} _ {2}, p _ {1} \times p _ {2})
\]

是A-模 \(\mathbf{M} \times \mathbf{M}\) 的分解，以及同态 \(\Delta : x \mapsto (x, x)\)（从 \(\mathbf{M}\) 到 \(\mathbf{M} \times \mathbf{M}\)）。根据引理 4，存在一个分解 \((\mathbf{Q}, q)\)（\(\mathbf{M}\) 的 \(\mathcal{C}_0\) 型有界分解）和一个交换图

\[
\begin{array}{c c c} \mathbf {Q} & \stackrel {{u}} {{\longrightarrow}} \mathbf {P} _ {1} & \times \mathbf {P} _ {2} \\ q \Big \downarrow & & \Big \downarrow p _ {1} \times p _ {2} \\ \mathbf {M} & \stackrel {{\Delta .}} {{\longrightarrow}} \mathbf {M} & \times \mathbf {M}; \end{array}
\]

我们推出一个交换图

\[
\begin{array}{c c c} \mathrm{Q} & \xrightarrow {u \circ p r _ {i}} & \mathrm{P} _ {i} \\ q \Big \downarrow & & \Big \downarrow p _ {i} \\ \mathbf {M} & \xrightarrow {\mathrm{l} _ {\mathrm{M}}} & \mathbf {M}, \end{array} \qquad i = 1, 2.
\]

根据引理 5，Con \((u \circ pr_i)\) 具有零同调，因此 \(u \circ pr_i\) 是一个拟同构且 \(\chi_{\mathcal{C}_0}(\mathbf{Q}) = \chi_{\mathcal{C}_0}(\mathbf{P}_i)\) (X，第 41 页，命题 10)；由此得出 \(\chi_{\mathcal{C}_0}(\mathbf{P}_1) = \chi_{\mathcal{C}_0}(\mathbf{P}_2)\) 如所宣称。

\begin{enumerate}
\def\labelenumi{\alph{enumi})}
\setcounter{enumi}{1}
\tightlist
\item
  对于 C 型的每个A-模 M，设 \(\varphi(M) \in K(\mathcal{C}_{0})\) 为 \(\chi_{\mathcal{C}_{0}}(P)\) 对 M 的所有 \(C_{0}\) 型有界左分解 P 所取的公共值。让我们证明函数 \(\varphi : \mathcal{C} \to K(\mathcal{C}_{0})\) 是可加的。为此设
\end{enumerate}

\[
0 \to \mathbf {M} ^ {\prime} \xrightarrow {f} \mathbf {M} \to \mathbf {M} ^ {\prime \prime} \to 0
\]

设为一个C型A-模的正合序列。根据引理4，存在一个交换图

\[
\begin{array}{c} \mathrm{P} ^ {\prime} \xrightarrow {u} \mathrm{P} \\ p ^ {\prime} \Big \downarrow \qquad \qquad \qquad \Big \downarrow p \\ 0 \longrightarrow \mathrm{M} ^ {\prime} \xrightarrow {f} \mathrm{M} \longrightarrow \mathrm{M} ^ {\prime \prime} \longrightarrow 0 \end{array}
\]

其中 (P, p) 和 (P', p') 是 \(\mathcal{C}_0\) 型的有界左分解。于是我们有

\[
\varphi (\mathbf {M}) = \chi_ {\mathcal {C} _ {0}} (\mathbf {P}), \quad \varphi (\mathbf {M} ^ {\prime}) = \chi_ {\mathcal {C} _ {0}} (\mathbf {P} ^ {\prime})
\]

并且由引理5

\[
\varphi (\mathbf {M} ^ {\prime \prime}) = \chi_ {\mathcal {C} _ {0}} (\operatorname{Con} (u)) = \chi_ {\mathcal {C} _ {0}} (\mathrm{P}) - \chi_ {\mathcal {C} _ {0}} (\mathrm{P} ^ {\prime}) = \varphi (\mathrm{M}) - \varphi (\mathrm{M} ^ {\prime});
\]

这正是我们想要证明的。

\begin{enumerate}
\def\labelenumi{\alph{enumi})}
\setcounter{enumi}{2}
\tightlist
\item
  设此时 \(\beta: \mathrm{K}(\mathcal{C}) \to \mathrm{K}(\mathcal{C}_0)\) 是这样一个同态，使得在上述记号下，我们有 \(\beta([\mathrm{M}]_{\mathcal{C}}) = \chi_{\mathcal{C}_0}(\mathrm{P})\) 。由于 \(p\) 是一个拟同构，我们有 \(\chi_{\mathcal{C}}(\mathrm{P}) = [\mathrm{M}]_{\mathcal{C}}\) ，因此 \(\alpha \circ \beta([\mathrm{M}]_{\mathcal{C}}) = \alpha(\chi_{\mathcal{C}_0}(\mathrm{P})) = \chi_{\mathcal{C}}(\mathrm{P}) = [\mathrm{M}]_{\mathcal{C}}\) ，故 \(\alpha \circ \beta = 1_{\mathrm{K}(\mathcal{C})}\) 。如果 M 是 \(\mathcal{C}_0\) 型的，那么 (M, \(1_{\mathrm{M}}\) ) 是 M 的一个分解，因此 \(\varphi(\mathrm{M}) = [\mathrm{M}]_{\mathcal{C}_0}\) ，且 \(\beta \circ \alpha = 1_{\mathrm{K}(\mathcal{C}_0)}\) 也成立，这就完成了证明。
\end{enumerate}

我们将在 § 8（X，第 137 页）中将此定理应用于“有限投射维数”的模。

\section{挠积}

在§§4至8中，我们用k表示一个交换环，用A表示一个结合的单位k-代数。k的作用本质上是辅助性的；我们主要考虑以下三种特殊情况：

\begin{enumerate}
\def\labelenumi{\alph{enumi})}
\item
  考虑一个任意环 A，设 k = Z，并赋予 A 其自然的 Z-代数结构，
\item
  考虑一个任意环 A，取 k 为 A 的中心，
\item
  考虑一个交换环 A，并取 k = A。
\end{enumerate}

\subsection{两个复形的张量积}

设 (C, d) 为右 A-模复形，(C', d') 为左 A-模复形。

赋予k-模 \(C \otimes_{A} C'\) 以如下分次，使得

\[
(\mathrm{C} \otimes_ {\mathbf {A}} \mathrm{C} ^ {\prime}) _ {n} = \sum_ {p + q = n} (\mathrm{C} _ {p} \otimes \mathrm{C} _ {q} ^ {\prime})
\]

并令D表示唯一的次数为 \((-1)\) 的k-线性自同态，它作用于 C \(\otimes_{A} C'\)，使得

\[
\mathbf {D} (x \otimes x ^ {\prime}) = d x \otimes x ^ {\prime} + (- 1) ^ {p} x \otimes d ^ {\prime} x ^ {\prime}, \quad x \in \mathrm{C} _ {p}, y \in \mathrm{C} _ {q} ^ {\prime}, p, q \in \mathbf {Z}. \tag {1}
\]

我们有 \(D \circ D = 0\)，因为依（1）中的记号，

\[
\mathrm{D} ^ {2} (x \otimes x ^ {\prime}) = d d x \otimes x ^ {\prime} + (- 1) ^ {p - 1} d x \otimes d ^ {\prime} x ^ {\prime} + (- 1) ^ {p} d x \otimes d ^ {\prime} x ^ {\prime} - x \otimes d ^ {\prime} d ^ {\prime} x ^ {\prime}.
\]

\(k\) -模复形 \((\mathbf{C} \otimes_{\mathbf{A}} \mathbf{C}', \mathbf{D})\) 称为复形 \((\mathbf{C}, d)\) 与 \((\mathbf{C}', d')\) 的张量积复形。

注记. --- 1) 当 \(\mathbf{C}'\) 简化为 \(\mathbf{C}_0' = \mathbf{M}\) 时，则 \((\mathbf{C} \otimes_{\mathbf{A}} \mathbf{C}')_n = \mathbf{C}_n \otimes_{\mathbf{A}} \mathbf{M}\) 且 \(D = d \otimes 1_M\)；例如 \(\mathbf{C} \otimes_{\mathbf{A}} \mathbf{A}_s\) 自然等同于 \(C\) 。类似地，当 \(C\) 简化为 \(\mathbf{C}_0 = P\) 时，则 \((\mathbf{C} \otimes_{\mathbf{A}} \mathbf{C}')_n = P \otimes_{\mathbf{A}} \mathbf{C}_n'\) 且 \(D = 1_P \otimes d\)。

\begin{enumerate}
\def\labelenumi{\arabic{enumi})}
\setcounter{enumi}{1}
\tightlist
\item
  对每个整数 \(r\)，我们有 \((C \otimes_{A} C')(r) = C(r) \otimes_{A} C'\) ，但 \((C \otimes_{A} C')(r)\) 与 \(C \otimes_{A} C'(r)\) 一般不具有相同的微分。
\end{enumerate}

设 \(p, q\) 是两个整数， \(x \in Z_p(C)\) , \(x' \in Z_q(C')\) ；则元素 \(x \otimes x'\) 作为 \(C_p \otimes C_q'\) 的元素依(1)属于 \(Z_{p+q}(C \otimes C')\)；此外，若 \(y \in C_{p+1}\) , \(y' \in C_{q+1}'\) ，我们有

\[
(x + d y) \otimes (x ^ {\prime} + d ^ {\prime} y ^ {\prime}) = x \otimes x ^ {\prime} + \mathrm{D} (y \otimes x ^ {\prime} + (- 1) ^ {p} (x + d y) \otimes y ^ {\prime})
\]

通过取商，我们导出一个k-线性映射，称为典范映射

\[
\gamma_ {p, q} (\mathrm{C}, \mathrm{C} ^ {\prime}): \mathrm{H} _ {p} (\mathrm{C}) \otimes_ {\mathrm{A}} \mathrm{H} _ {q} (\mathrm{C} ^ {\prime}) \rightarrow \mathrm{H} _ {p + q} (\mathrm{C} \otimes_ {\mathrm{A}} \mathrm{C} ^ {\prime});
\]

如果我们装备 \(\mathbf{H}(\mathbf{C})\otimes_{\mathbf{A}}\mathbf{H}(\mathbf{C}^{\prime})\) 以如下分次，使得

\[
\left(\mathrm{H} (\mathrm{C}) \otimes \mathrm{H} (\mathrm{C} ^ {\prime})\right) _ {n} = \sum_ {p + q = n} \mathrm{H} _ {p} (\mathrm{C}) \otimes_ {\mathrm{A}} \mathrm{H} _ {p} (\mathrm{C} ^ {\prime}),
\]

\(\gamma_{p,q}\) 定义一个次数为0的分次 \(k\) -线性映射

\[
\gamma (\mathbf {C}, \mathbf {C} ^ {\prime}): \mathrm{H} (\mathbf {C}) \otimes_ {\mathbf {A}} \mathrm{H} (\mathbf {C} ^ {\prime}) \rightarrow \mathrm{H} (\mathbf {C} \otimes_ {\mathbf {A}} \mathbf {C} ^ {\prime}).
\]

II，第59页的命题6改述如下：

如果复形C和C'在右侧为零，则C ⊗\_A C' 在右侧为零，且典范的 k-线性映射

\[
\gamma_ {0, 0} (\mathbf {C}, \mathbf {C} ^ {\prime}): \mathrm{H} _ {0} (\mathbf {C}) \otimes_ {\mathbf {A}} \mathrm{H} _ {0} (\mathbf {C} ^ {\prime}) \rightarrow \mathrm{H} _ {0} (\mathbf {C} \otimes_ {\mathbf {A}} \mathbf {C} ^ {\prime})
\]

是双射。

设 \(u: (\mathbf{C}, d) \to (\mathbf{C}_1, d_1)\) 是右A-模复形的一个态射，且 \(u': (\mathbf{C}', d') \to (\mathbf{C}_1', d_1')\) 是左A-模复形的一个态射；则 \(u \otimes u': \mathbf{C} \otimes_{\mathbf{A}} \mathbf{C}' \to \mathbf{C}_1 \otimes_{\mathbf{A}} \mathbf{C}_1'\) 是 \(k\) -模复形的态射；事实上，它是次数为0的分次态射，并且若用D和 \(\mathbf{D}_1\) 表示 \(\mathbf{C} \otimes \mathbf{C}'\) 与 \(\mathbf{C}_1 \otimes \mathbf{C}_1'\) 的微分，则对于 \(p, q \in \mathbf{Z}, x \in \mathbf{C}_p, x' \in \mathbf{C}_q\) 我们有，

\[
\begin{array}{c} (u \otimes u ^ {\prime}) (\mathrm{D} (x \otimes x ^ {\prime})) = u (d x) \otimes u ^ {\prime} (x ^ {\prime}) + (- 1) ^ {p} u (x) \otimes u ^ {\prime} (d ^ {\prime} x ^ {\prime}) = \\ = d _ {1} u (x) \otimes u ^ {\prime} (x ^ {\prime}) + (- 1) ^ {p} u (x) \otimes d _ {1} ^ {\prime} u ^ {\prime} (x ^ {\prime}) = \mathrm{D} _ {1} (u (x) \otimes u ^ {\prime} (x ^ {\prime}))  . \end{array}
\]

此外，以下图表是可交换的：

\[
\begin{array}{c} \mathrm{H(C)} \otimes_ {\mathrm{A}} \mathrm {H(C^ {\prime})} \xrightarrow {\gamma (\mathrm{C,C} ^ {\prime})} \mathrm{H(C} \otimes_ {\mathrm{A}} \mathrm {C^ {\prime}}) \\ \mathrm{H(u)} \otimes \mathrm {H(u^{\prime})} \Big \downarrow \\ \mathrm {H(C_ {1})} \otimes_ {\mathrm{A}} \mathrm {H(C_ {1} ^ {\prime})} \xrightarrow [ \gamma (\mathrm {C_ {1}} , \mathrm {C_ {1} ^ {\prime}}) ]{} \mathrm {H(C_ {1}} \otimes_ {\mathrm{A}} \mathrm {C_ {1} ^ {\prime}}). \end{array}
\]

设A°为A的反k-代数，并设C°（相应地C'°）为将复形C（相应地C'）视为左（相应地右）A°-模的复形。

\[
\sigma (\mathbf {C}, \mathbf {C} ^ {\prime}): \mathbf {C} \otimes_ {\mathbf {A}} \mathbf {C} ^ {\prime} \rightarrow \mathbf {C} ^ {\prime \circ} \otimes_ {\mathbf {A} ^ {\circ}} \mathbf {C} ^ {\circ}
\]

唯一的次数为0的分次 \(k\) -线性映射，使得对于 \(x \in C_p\) , \(x' \in C_q'\) , \(p, q \in \mathbf{Z}\) ，有

\[
\sigma (\mathbf {C}, \mathbf {C} ^ {\prime}) (x \otimes x ^ {\prime}) = (- 1) ^ {p q} x ^ {\prime} \otimes x.
\]

命题2. — 映射 \(\sigma(\mathbf{C},\mathbf{C}'):\mathbf{C}\otimes_{\mathbf{A}}\mathbf{C}'\to\mathbf{C}^{\prime\circ}\otimes_{\mathbf{A}^{\circ}}\mathbf{C}^{\circ}\) 是 \(k\) -模复形的同构，其逆同构为 \(\sigma(\mathbf{C}^{\prime\circ},\mathbf{C}^{\circ})\) 。

由于映射 \(\sigma(\mathbf{C},\mathbf{C}^{\prime})\) 与 \(\sigma(\mathbf{C}^{\circ},\mathbf{C}^{\circ})\) 互为逆映射，只需证明 \(\sigma(\mathbf{C},\mathbf{C}^{\prime})\) 是复形的态射。现在，对于

\[
x \in \mathrm{C} _ {p} = \mathrm{C} _ {p} ^ {\circ}, \quad x ^ {\prime} \in \mathrm{C} _ {q} ^ {\prime} = \mathrm{C} _ {q} ^ {\prime \circ}, \quad p, q \in \mathbf {Z},
\]

我们有，用 \(D^{c}\) 记 \(C' \otimes_{A^{0}} C\) 的微分，

\[
\begin{array}{l} \sigma (\mathrm{C}, \mathrm{C} ^ {\prime}) \circ \mathrm{D} (x \otimes x ^ {\prime}) = \sigma (\mathrm{C}, \mathrm{C} ^ {\prime}) (d x \otimes x ^ {\prime} + (- 1) ^ {p} x \otimes d ^ {\prime} x ^ {\prime}) = \\ = (- 1) ^ {(p + 1) q} x ^ {\prime} \otimes d x + (- 1) ^ {p + p (q + 1)} d ^ {\prime} x ^ {\prime} \otimes x = (- 1) ^ {p q} d ^ {\prime} x ^ {\prime} \otimes x + (- 1) ^ {p q + q} x ^ {\prime} \otimes d x \\ = (- 1) ^ {p q} \mathrm{D} ^ {\circ} (x ^ {\prime} \otimes x) = \mathrm{D} ^ {\circ} \circ \sigma (\mathrm{C}, \mathrm{C} ^ {\prime}) (x \otimes x ^ {\prime}); \end{array}
\]

这给出 \(\sigma(C, C') \circ D = D^{\circ} \circ \sigma(C, C')\) ，由此得到所期望的断言。

同构 \(\sigma(\mathbf{C},\mathbf{C}^{\prime}):\mathbf{C}\otimes_{\mathbf{A}}\mathbf{C}^{\prime}\to\mathbf{C}^{\circ}\otimes_{\mathbf{A}^{\circ}}\mathbf{C}^{-}\) 称为复形 \(\mathbf{C}\) 与 \(\mathbf{C}^{\prime}\) 的张量积的交换同构。

若 \(u: \mathbf{C} \to \mathbf{C}_{1}\) 与 \(v: \mathbf{C}' \to \mathbf{C}_{1}'\) 是如上所述的两个复形的态射，我们有如下交换图：

\[
\begin{array}{c} \mathrm{C} \otimes_ {\mathrm{A}} \mathrm{C} ^ {\prime} \xrightarrow {\sigma (\mathrm{C} , \mathrm{C} ^ {\prime})} \mathrm{C} ^ {\prime \circ} \otimes_ {\mathrm{A} ^ {\circ}} \mathrm{C} ^ {\circ} \\ u \otimes u ^ {\prime} \Bigg \downarrow \\ \mathrm{C} _ {1} \otimes_ {\mathrm{A}} \mathrm{C} _ {1} ^ {\prime} \xrightarrow [ \sigma (\mathrm{C} _ {1} , \mathrm{C} _ {1} ^ {\prime}) ]{} \mathrm{C} _ {1} ^ {\prime \circ} \otimes_ {\mathrm{A} ^ {\circ}} \mathrm{C} _ {1} ^ {\circ}. \end{array}
\]

假设在本节剩余部分中环 A 是交换的（一般情形参见第 9 节）。

设 C, C′, C′′ 为三个 A-模复形；A-模的典范同态（III，第 64 页）

\[
\varphi : \left(\mathrm{C} \otimes_ {\mathrm{A}} \mathrm{C} ^ {\prime}\right) \otimes_ {\mathrm{A}} \mathrm{C} ^ {\prime \prime} \rightarrow \mathrm{C} \otimes_ {\mathrm{A}} \left(\mathrm{C} ^ {\prime} \otimes_ {\mathrm{A}} \mathrm{C} ^ {\prime \prime}\right)
\]

是复形的同构，这一点可直接由定义验证。

更一般地，设 \((\mathbf{C}^{(i)}, d^{(i)})_{i \in I}\) 是一族A-模复形，其中指标集I是有限且全序的；为记号简便，我们将I等同于N中的区间 \([1, r]\) 。赋予A-模 \(C = \bigotimes_{i=1}^{r} C^{(i)}\) 以如下分次，使得

\[
\mathrm{C} _ {n} = \sum_ {p _ {1} + p _ {2} + \dots + p _ {r} = n} (\mathrm{C} ^ {(1)}) _ {p _ {1}} \otimes (\mathrm{C} ^ {(2)}) _ {p _ {2}} \otimes \dots \otimes (\mathrm{C} ^ {(r)}) _ {p _ {r}},
\]

，并定义C上的一个次数为(-1)的分次A-自同态为

\[
\mathrm{D} \left(x _ {1} \otimes \dots \otimes x _ {r}\right) = \sum_ {j = 1} ^ {r} (- 1) ^ {p _ {1} + \dots + p _ {j - 1}} x _ {1} \otimes \dots \otimes x _ {j - 1} \otimes d _ {j} x _ {j} \otimes x _ {j + 1} \otimes \dots \otimes x _ {r}
\]

其中 \(x_{i}\in (\mathbf{C}^{(i)})_{p_{i}}\)，\(i = 1,\dots,n\)。则(C, D)是一个A-模复形，称为族 \((\mathbf{C}_{i}, d_{i})\) 的张量积复形。对 [0, r] 中的每一个严格递增序列 \(r_{0}, \ldots, r_{k}\)，其中 \(r_{0} = 0\) , \(r_{k} = r\) ，典范结合性同构

\[
\bigotimes_ {j = 0} ^ {k - 1} \left(\bigotimes_ {i = r _ {j} + 1} ^ {r _ {j + 1}} \mathbf {C} ^ {(i)}\right)\rightarrow \bigotimes_ {i = 1} ^ {r} \mathbf {C} ^ {(i)}
\]

是复形的同构。

我们如上定义次数为 0 的分次同态

\[
\gamma ((\mathbf {C} ^ {(i)})): \underset {i \in \mathbf {I}} {\otimes} \mathrm{H} (\mathbf {C} ^ {(i)}) \to \mathrm{H} \left(\underset {i \in \mathbf {I}} {\otimes} \mathbf {C} ^ {(i)}\right) .
\]

注记。--- 3) 可以定义有限族复形的张量积，而无需赋予指标集全序（X，第 185 页，练习 3）。

\begin{enumerate}
\def\labelenumi{\arabic{enumi})}
\setcounter{enumi}{3}
\tightlist
\item
  假设每个 \(\mathbf{C}^{(i)}\) 被赋予一个与其分次结构相容的分次代数结构，且使得 \(d^{(i)}\) 是反导子（III，第117页）。然后让我们赋予 \(\bigotimes_{i\in I}\mathbf{C}^{(i)}\) 以给定结构的左分次张量积代数结构（III，第49页）。则D是一个反导子。事实上，利用张量积的结合性，可以假设I = \{1, 2\}；然后设 \(p_1\) , \(q_1\) , \(p_2\) , \(q_2 \in \mathbf{Z}\) , \(x_1 \in (\mathbf{C}^{(1)})_{p_1}\) , \(y_1 \in (\mathbf{C}^{(1)})_{q_1}\) , \(x_2 \in (\mathbf{C}^{(2)})_{p_2}\) , \(y_2 \in (\mathbf{C}^{(2)})_{q_2}\) ；我们有
\end{enumerate}

\[
\begin{array}{l} \left(\mathrm{D} (x _ {1} \otimes x _ {2})\right) (y _ {1} \otimes y _ {2}) + (- 1) ^ {p _ {1} + p _ {2}} (x _ {1} \otimes x _ {2}) (\mathrm{D} (y _ {1} \otimes y _ {2})) = \\ = (d x _ {1} \otimes x _ {2} + (- 1) ^ {p _ {1}} x _ {1} \otimes d x _ {2}) (y _ {1} \otimes y _ {2}) + \\ \quad + (- 1) ^ {p _ {1} + p _ {2}} (x _ {1} \otimes x _ {2}) (d y _ {1} \otimes y _ {2} + (- 1) ^ {q _ {1}} y _ {1} \otimes d y _ {2}) = \\ = (- 1) ^ {p _ {2} q _ {1}} (d x _ {1}) y _ {1} \otimes x _ {2} y _ {2} + (- 1) ^ {p _ {1} + (p _ {2} - 1) q _ {1}} x _ {1} y _ {1} \otimes (d x _ {2}) y _ {2} + \\ \quad + (- 1) ^ {p _ {1} + p _ {2} + p _ {2} (q _ {1} - 1)} x _ {1} d y _ {1} \otimes x _ {2} y _ {2} + (- 1) ^ {p _ {1} + p _ {2} + q _ {1} + p _ {2} q _ {1}} x _ {1} y _ {1} \otimes x _ {2} d y _ {2} \\ = (- 1) ^ {p _ {2} q _ {1}} [ (d x _ {1}) y _ {1} + (- 1) ^ {p _ {1}} x _ {1} d y _ {1} ] \otimes x _ {2} y _ {2} + \\ \quad + (- 1) ^ {p _ {1} + q _ {1} + p _ {2} q _ {1}} x _ {1} y _ {1} \otimes ((d x _ {2}) y _ {2} + (- 1) ^ {p _ {2}} x _ {2} d y _ {2}) \\ = (- 1) ^ {p _ {2} q _ {1}} [ d (x _ {1} y _ {1}) \otimes x _ {2} y _ {2} + (- 1) ^ {p _ {1} + q _ {1}} x _ {1} y _ {1} \otimes d (x _ {2} y _ {2}) ] \\ = (- 1) ^ {p _ {2} q _ {1}} \mathrm{D} (x _ {1} y _ {1} \otimes x _ {2} y _ {2}) = \mathrm{D} ((x _ {1} \otimes x _ {2}) (y _ {1} \otimes y _ {2}))  . \end{array}
\]

\subsection{张量积与同伦}

命题3. --- 设C, C \(_{1}\) 为右A-模的两个复形，C'，C' \(_{1}'\) 为左A-模的两个复形，以及u：C → C \(_{1}\) ，v：C → C \(_{1}\) ，u' : C' → C \(_{1}'\) ，v' : C' → C \(_{1}'\) 是复形的态射。

\begin{enumerate}
\def\labelenumi{\alph{enumi})}
\item
  若 \(u\) 与 \(u'\) 分别同伦于 \(v\) 与 \(v'\)，则两个态射 \(u \otimes u'\) 与 \(v \otimes v'\)（从 \(\mathbf{C} \otimes_{\mathbf{A}} \mathbf{C}'\) 到 \(\mathbf{C}_1 \otimes_{\mathbf{A}} \mathbf{C}_1'\)）是同伦的。
\item
  如果 u 和 \(u'\) 是同伦等价， \(u \otimes u'\) 是一个同伦等价。
\item
  如果 C 或 \(\mathbf{C}'\) 同伦于零，则 \(\mathbf{C} \otimes_{\mathbf{A}} \mathbf{C}'\) 同伦于零。
\end{enumerate}

让我们用同一个字母 d 表示复形 C, \(C_{1}\), \(C'\), \(C_{1}'\) 的微分，并用 D 表示复形 \(C \otimes_{A} C'\) 与 \(C_{1} \otimes_{A} C_{1}'\) 的微分。

若 \(u\) (相应地， \(u'\) ) 同伦于 \(v\) (相应地， \(v'\) )，则存在一个次数为 1 的分次同态 \(s: \mathbf{C} \to \mathbf{C}_1\) (相应地， \(s': \mathbf{C}' \to \mathbf{C}_1'\) ）使得

\[
u - v = d s + s d \qquad (\mathrm{resp.} u ^ {\prime} - v ^ {\prime} = d s ^ {\prime} + s ^ {\prime} d).\tag{2}
\]

设 \(\mathbf{S}:\mathbf{C}\otimes_{\mathbf{A}}\mathbf{C}^{\prime}\to \mathbf{C}_{1}\otimes_{\mathbf{A}}\mathbf{C}_{1}^{\prime}\) 是唯一的次数为1的分次同态，使得对于 \(x\in \mathbf{C}_p\) , \(y\in \mathbf{C}_q^\prime\) , \(p,q\in \mathbf{Z}\) ，有

\[
\mathbf {S} (x \otimes y) = s (x) \otimes u ^ {\prime} (y) + (- 1) ^ {p} v (x) \otimes s ^ {\prime} (y).\tag{3}
\]

于是我们有，采用前述记号：

\[
\begin{array}{r l} (\mathrm{DS} + \mathrm{SD}) (x \otimes y) & = \mathrm{D} (s x \otimes u ^ {\prime} y) + (- 1) ^ {p} \mathrm{D} (v x \otimes s ^ {\prime} y) + \mathrm{S} (d x \otimes y) + \\ & + (- 1) ^ {p} \mathrm{S} (x \otimes d y) = \\ & = d s x \otimes u ^ {\prime} y + (- 1) ^ {p + 1} s x \otimes d u ^ {\prime} y + (- 1) ^ {p} d v x \otimes s ^ {\prime} y + v x \otimes d s ^ {\prime} y + \\ & + s d x \otimes u ^ {\prime} y + (- 1) ^ {p - 1} v d x \otimes s ^ {\prime} y + (- 1) ^ {p} s x \otimes u ^ {\prime} d y + v x \otimes s ^ {\prime} d y \\ & = (d s + s d) (x) \otimes u ^ {\prime} y + v x \otimes (d s ^ {\prime} + s ^ {\prime} d) (y) \\ & = (u x - v x) \otimes u ^ {\prime} y + v x \otimes (u ^ {\prime} y - v ^ {\prime} y) = u x \otimes u ^ {\prime} y - v x \otimes v ^ {\prime} y. \end{array}
\]

这给出 DS + SD = u ⊗ u' - v ⊗ v'，由此得 a)。

让我们证明b)。如果u和 \(u'\) 是同伦等价，则存在复形的同态 \(\alpha: C_{1} \to C\) 与 \(\alpha': C_{1}' \to C'\) 使得 \(u \circ \alpha\) , \(\alpha \circ u\) , \(u' \circ \alpha'\) , \(\alpha' \circ u'\) 分别同伦于 \(Id_{C_{1}}\) , \(Id_{C}\) , \(Id_{C_{1}}\) 与 \(Id_{C'}\) 。而 \((u \otimes u') \circ (\alpha \otimes \alpha')\) 等于 \((u \circ \alpha) \otimes (u' \circ \alpha')\) ，由a)它同伦于 \(Id_{C_{1}} \otimes Id_{C_{1}} = Id_{C_{1} \otimes C_{1}}\) ；而 \((\alpha \otimes \alpha') \circ (u \otimes u')\) 同伦于 \(Id_{C \otimes C'}\) ，由此得 b)。最后，c) 由 b) 应用于 \(C_{1}\) 或 \(C_{1}'\) 为零的情形得出。

推论 1. --- 设 C' 是左 A-模的一个分裂复形，使得 H(C') 是平坦的。对于右 A-模的每个复形 C，典范映射

\[
\gamma (\mathrm{C}, \mathrm{C} ^ {\prime}): \mathrm{H} (\mathrm{C}) \otimes_ {\mathrm{A}} \mathrm{H} (\mathrm{C} ^ {\prime}) \rightarrow \mathrm{H} (\mathrm{C} \otimes_ {\mathrm{A}} \mathrm{C} ^ {\prime})
\]

是双射。

根据X，第35页，定义6，存在一个同伦等价 \(u': C' \to H(C')\) 。由命题3， \(1_{C} \otimes u': C \otimes_{A} C' \to C \otimes_{A} H(C')\) 是一个同伦等价；由于

\[
\mathrm{H} \left(1 _ {\mathrm{C}} \otimes u ^ {\prime}\right) \circ \gamma (\mathrm{C}, \mathrm{C} ^ {\prime}) = \gamma (\mathrm{C}, \mathrm{H} (\mathrm{C} ^ {\prime})) \circ \left(1 _ {\mathrm{C}} \otimes \mathrm{H} (u ^ {\prime})\right),
\]

由于 \(\mathbf{H}(1_{\mathbf{C}} \otimes u')\) 与 \(\mathbf{H}(u')\) 都是双射的，只需证明 \(\gamma(\mathbf{C}, \mathbf{H}(\mathbf{C}'))\) 是双射的，于是我们归结为 \(\mathbf{C}'\) 是平坦且微分为零的情形。在此情形下，典范正合序列

\begin{enumerate}
\def\labelenumi{(\Roman{enumi})}
\tightlist
\item
\end{enumerate}

\[
0 \to \mathbf {Z} (\mathbf {C}) \xrightarrow {i} \mathbf {C} \xrightarrow {\delta} \mathbf {B} (\mathbf {C}) \to 0\tag{II}
\]

\[
0 \rightarrow \mathrm{B} (\mathrm{C}) \xrightarrow {j} \mathrm{Z} (\mathrm{C}) \xrightarrow {\pi} \mathrm{H} (\mathrm{C}) \rightarrow 0
\]

给出正合序列：

\[
\begin{array}{l} 0 \longrightarrow Z (C) \otimes_ {A} C ^ {\prime} \xrightarrow {i \otimes 1} C \otimes_ {A} C ^ {\prime} \quad \xrightarrow {\delta \otimes 1} B (C) \otimes_ {A} C ^ {\prime} \longrightarrow 0 \\ 0 \longrightarrow B (C) \otimes_ {A} C ^ {\prime} \xrightarrow {j \otimes 1} Z (C) \otimes_ {A} C ^ {\prime} \xrightarrow {\pi \otimes 1} H (C) \otimes_ {A} C ^ {\prime} \longrightarrow 0. \end{array}
\]

由于 \(d = i \circ j \circ \delta\)，我们有 \(\mathbf{D} = d \otimes 1_{\mathbf{C}'} = (i \otimes 1) \circ (j \otimes 1) \circ (\delta \otimes 1)\)，这表明典范映射 \(\mathbf{Z}(\mathbf{C}) \otimes_{\mathbf{A}} \mathbf{C}' \to \mathbf{Z}(\mathbf{C} \otimes_{\mathbf{A}} \mathbf{C}')\) 与 \(\mathbf{B}(\mathbf{C}) \otimes_{\mathbf{A}} \mathbf{C}' \to \mathbf{B}(\mathbf{C} \otimes_{\mathbf{A}} \mathbf{C}')\) 是双射，因此通过取商，\(\gamma(\mathbf{C}, \mathbf{C}')\) 也是双射的。

推论2. --- 设N为平坦的左A-模。对于右A-模的每个复形C，典范同态

\[
\gamma_ {n} (\mathbf {C}, \mathbf {N}): \mathrm{H} _ {n} (\mathbf {C}) \otimes_ {\mathbf {A}} \mathbf {N} \rightarrow \mathrm{H} _ {n} (\mathbf {C} \otimes_ {\mathbf {A}} \mathbf {N})
\]

都是双射的。

推论3. --- 设C'为左A-模的复形，使得B(C')和H(C')是投射的。对于右A-模的每个复形C，映射γ(C, C')是双射。

事实上，C'是分裂的（X，第35页，例4），且H(C')是投射的；因此可以应用推论1。

注记. --- 1) 利用交换同构，由推论1、2和3可推出将张量积的两个参数角色互换后所得的类似结论。

\begin{enumerate}
\def\labelenumi{\arabic{enumi})}
\setcounter{enumi}{1}
\tightlist
\item
  我们将在下文（X，第79页，推论4）看到，当假设C'和H(C')是平坦的且C'在右侧有界时，推论1的结论也成立。
\end{enumerate}

\subsection{与一个右有界平坦复形作张量积}

引理1. --- 设C为右A-模复形，E为左A-模复形。假设H(C) = 0且E是平坦的且在右侧有界。则H(C ⊗\_A E) = 0。

对于 \(k \in \mathbf{Z}\) ，设 \(\mathbf{T}^{(k)}\) 为 \(\mathbf{C} \otimes_{\mathbf{A}} \mathbf{E}\) 的子复形，使得

\[
\mathrm{T}_{n}^{(k)} = \sum_{\substack{p + q = n\\ q\leqslant k}}\mathrm{C}_{p}\otimes_{\mathrm{A}}\mathrm{E}_{q};
\]

那么 \(\mathbf{T}^{(k - 1)}\subset \mathbf{T}^{(k)}\) 并且我们有一个复形的正合序列

\[
0 \longrightarrow \mathrm{T} ^ {(k - 1)} \stackrel {i _ {k}} {\longrightarrow} \mathrm{T} ^ {(k)} \stackrel {\pi} {\longrightarrow} \mathrm{C} \otimes_ {\mathbf {A}} \mathrm{E} _ {k} (- k) \longrightarrow 0
\]

其中 \(i_{k}\) 是典范单射，\(\pi\) 将前面的直和投影到其因子 \(\mathbf{C}_{n-k} \otimes_{\mathbf{A}} \mathbf{E}_{k} = (\mathbf{C} \otimes_{\mathbf{A}} \mathbf{E}_{k}(-k))_{n}\) 上。由上述推论2，我们有 \(\mathbf{H}(\mathbf{C} \otimes_{\mathbf{A}} \mathbf{E}_{k}(-k)) = 0\) ，因此 \(i_{k}\) 是一个拟同构。我们有 \(\mathbf{T}^{(k)} = 0\)（当 \(k\) 足够小时），因为E在右侧有界，因此对k归纳可得对所有k有 \(\mathrm{H}(\mathrm{T}^{(k)}) = 0\)。最后，典范态射 \(\lim_{\mathbf{A}} \mathrm{T}^{(k)} \to \mathbf{C} \otimes_{\mathbf{A}} \mathbf{E}\) 显然是一个同构，因此 \(\mathrm{H}(\mathbf{C} \otimes_{\mathbf{A}} \mathbf{E}) = 0\) (X，第28页，命题1)。

引理2. --- 若 \(u: \mathbf{C} \to \mathbf{C}'\) 是右A-模复形的一个态射，且E是左A-模复形，则复形Con \((u) \otimes_{\mathbf{A}} \mathbf{E}\) 以及 Con \((u \otimes 1_{\mathbf{E}})\) 是同构的。

根据定义，Con \((u) \otimes_{\mathbf{A}} \mathbf{E}\) 是分次模 \((\mathbf{C}'(-1) \oplus \mathbf{C}) \otimes_{\mathbf{A}} \mathbf{E}\) 赋予微分D，使得对于 \(x \in \mathbf{C}_p\) , \(y' \in \mathbf{C}'(-1)_p = \mathbf{C}_{p-1}'\) , \(z \in \mathbf{E}_q\) ，有

\[
\mathrm{D} ((y ^ {\prime}, x) \otimes z) = (- d y ^ {\prime}, d x - u (y ^ {\prime})) \otimes z + (- 1) ^ {p} (y ^ {\prime}, x) \otimes d z,\tag{4}
\]

而 \(\operatorname{Con}(u \otimes 1_{\mathbf{E}})\) 是分次模 \((\mathbf{C}' \otimes_{\mathbf{A}} \mathbf{E})(-1) \oplus (\mathbf{C} \otimes_{\mathbf{A}} \mathbf{E})\) 配备微分 \(\mathbf{D}_1\) 使得，对于 \(x \in \mathbf{C}_p\) , \(y' \in \mathbf{C}_{p-1}'\) , \(z \in \mathbf{E}_q\) ，有

\[
\mathrm{D} _ {1} \left(y ^ {\prime} \otimes z, x \otimes z\right) = (- d y ^ {\prime} \otimes z - (- 1) ^ {p - 1} y ^ {\prime} \otimes d z, d x \otimes z + (- 1) ^ {p} x \otimes d z - u \left(y ^ {\prime}\right) \otimes z)
\]

\[
= \left(- d y ^ {\prime} \otimes z, (d x - u (y ^ {\prime})) \otimes z\right) + (- 1) ^ {p} \left(y ^ {\prime} \otimes d z, x \otimes d z\right)
\]

由此得出该论断。

命题4. — 设 \(u: \mathbf{C} \to \mathbf{C}'\) 是右A-模复形间的一个拟同构，且E是左A-模的一个复形，平坦且在右侧有界。则

\[
u \otimes 1 _ {\mathbf {E}}: \mathbf {C} \otimes_ {\mathbf {A}} \mathbf {E} \rightarrow \mathbf {C} ^ {\prime} \otimes_ {\mathbf {A}} \mathbf {E}
\]

是k-模复形间的一个拟同构。

事实上，根据X，第38页，推论， \(u\) (相应地， \(u \otimes 1_{\mathrm{E}}\) ）是拟同构当且仅当 \(\mathrm{H}(\mathrm{Con}(u)) = 0\) (相应地， \(\mathrm{H}(\mathrm{Con}(u \otimes 1_{\mathrm{E}})) = 0\) ）。然后我们由引理1和引理2得出结论。

注记。--- 利用交换同构，由前面的陈述可推出将张量积的两个参数角色互换后所得的类似陈述。

\subsection{挠积的定义与基本性质}

对于每个A-模E，我们用 \(p_{E}: L(E) \to E\) 表示E的典范自由分解（X，第 50 页）。

定义1.——设M为右A-模，N为左A-模。M与N的挠积是分次k-模

\[
\operatorname{Tor} ^ {\mathbf {A}} (\mathbf {M}, \mathbf {N}) = \mathrm{H} (\mathrm{L} (\mathbf {M}) \otimes_ {\mathbf {A}} \mathrm{L} (\mathbf {N}))  .\tag{4}
\]

Tor \(^{A}\) (M, N) 的齐次分量记为

\[
\operatorname{Tor} _ {n} ^ {\mathbf {A}} (\mathbf {M}, \mathbf {N}) = \mathrm{H} _ {n} (\mathrm{L} (\mathbf {M}) \otimes_ {\mathbf {A}} \mathrm{L} (\mathbf {N})).\tag{5}
\]

由于L(M)和L(N)在右侧为零，我们有

\[
\operatorname{Tor} _ {n} ^ {\mathbf {A}} (\mathbf {M}, \mathbf {N}) = 0 \quad \text { for } \quad n <   0.\tag{6}
\]

注记1. --- 我们将在下文（X，第107页，命题6）看到Tor \(^{A}\) (M, N) 模的有限性性质。例如，若 A 是交换诺特环，且 M 与 N 是有限生成的 A-模，则每个 A-模 Tor \(_{n}^{A}\) (M, N) 是有限生成的。

设 \(f: \mathbf{M} \to \mathbf{M}'\) 是右A-模的同态，且 \(g: \mathbf{N} \to \mathbf{N}'\) 为左A-模的一个同态；我们令 \(\operatorname{Tor}^{\mathbf{A}}(f, g) = \operatorname{H}\left(\operatorname{L}(f) \otimes_{\mathbf{A}} \operatorname{L}(g)\right)\) ；它是分次 \(k\) -模的一个同态

\[
\operatorname{Tor} ^ {\mathbf {A}} (f, g): \operatorname{Tor} ^ {\mathbf {A}} (\mathbf {M}, \mathbf {N}) \rightarrow \operatorname{Tor} ^ {\mathbf {A}} (\mathbf {M} ^ {\prime}, \mathbf {N} ^ {\prime})
\]

其齐次分量记为

\[
\operatorname{Tor} _ {n} ^ {\mathbf {A}} (f, g): \operatorname{Tor} _ {n} ^ {\mathbf {A}} (\mathbf {M}, \mathbf {N}) \rightarrow \operatorname{Tor} _ {n} ^ {\mathbf {A}} (\mathbf {M} ^ {\prime}, \mathbf {N} ^ {\prime}).
\]

根据X，第62页，命题1，典范同态

\[
\gamma_ {0, 0}: \mathrm{H} _ {0} (\mathrm{L} (\mathrm{M})) \otimes_ {\mathrm{A}} \mathrm{H} _ {0} (\mathrm{L} (\mathrm{N})) \rightarrow \mathrm{H} _ {0} (\mathrm{L} (\mathrm{M}) \otimes_ {\mathrm{A}} \mathrm{L} (\mathrm{N}))
\]

是双射的；利用同构 \(\mathbf{M} \to \mathbf{H}_{0}(\mathbf{L}(\mathbf{M}))\) 与 \(\mathbf{N} \to \mathbf{H}_{0}(\mathbf{L}(\mathbf{N}))\)，我们推导出一个同构，称为典范同构

\[
\gamma_ {\mathbf {M}, \mathbf {N}}: \mathbf {M} \otimes_ {\mathbf {A}} \mathbf {N} \rightarrow \operatorname{Tor} _ {0} ^ {\mathbf {A}} (\mathbf {M}, \mathbf {N}).\tag{7}
\]

我们将始终通过这个同构把 \(\operatorname{Tor}_{0}^{\mathbf{A}}(\mathbf{M},\mathbf{N})\) 与 \(\mathbf{M}\otimes_{\mathbf{A}}\mathbf{N}\) 等同起来。于是 \(k\) -线性映射 \(\operatorname{Tor}_{0}^{\mathbf{A}}(f,g)\) 与 \(f\otimes g\) 等同。

注记2. --- 复形的态射 \(p_{\mathbf{M}} \otimes p_{\mathbf{N}}: \mathbf{L}(\mathbf{M}) \otimes_{\mathbf{A}} \mathbf{L}(\mathbf{N}) \to \mathbf{M} \otimes_{\mathbf{A}} \mathbf{N}\) 在0次同调上诱导同构

\[
\gamma_ {\mathbf {M}, \mathbf {N}} ^ {- 1}: \operatorname{Tor} _ {0} ^ {\mathbf {A}} (\mathbf {M}, \mathbf {N}) \rightarrow \mathbf {M} \otimes_ {\mathbf {A}} \mathbf {N}
\]

\(\gamma_{M,N}\) 的逆。

我们有 \(L(1_{\mathbf{M}}) = 1_{L(\mathbf{M})}\) , \(L(1_{\mathbf{N}}) = i_{L(\mathbf{N})}\) 因此通过取同调：

\[
\operatorname{Tor} ^ {\mathbf {A}} \left(1 _ {\mathbf {M}}, 1 _ {\mathbf {N}}\right) = 1 _ {\operatorname{Tor} ^ {\mathbf {A}} (\mathbf {M}, \mathbf {N})}.\tag{8}
\]

若 \(f': \mathbf{M}' \to \mathbf{M}''\) (相应地， \(g': \mathbf{N}' \to \mathbf{N}''\) )是右（相应地，左）A-模的同态，我们有 \(\mathbf{L}(g' \circ g) = \mathbf{L}(g') \circ \mathbf{L}(g)\) 与 \(\mathbf{L}(f' \circ f) = \mathbf{L}(f') \circ \mathbf{L}(f)\) ，因此

\[
\operatorname{Tor} ^ {\mathbf {A}} \left(f ^ {\prime} \circ f, g ^ {\prime} \circ g\right) = \operatorname{Tor} ^ {\mathbf {A}} \left(f ^ {\prime}, g ^ {\prime}\right) \circ \operatorname{Tor} ^ {\mathbf {A}} (f, g).\tag{9}
\]

考虑k-复形的态射

\[
\mathrm{L} (\mathrm{M}) \otimes_ {\mathrm{A}} \mathrm{N} \xleftarrow {1 \otimes p _ {\mathrm{N}}} \mathrm{L} (\mathrm{M}) \otimes_ {\mathrm{A}} \mathrm{L} (\mathrm{N}) \xrightarrow {p _ {\mathrm{M}} \otimes 1} \mathrm{M} \otimes_ {\mathrm{A}} \mathrm{L} (\mathrm{N})
\]

以及它们在同调中诱导的k-同态：

\[
\mathrm{H} (\mathrm{L} (\mathrm{M}) \otimes_ {\mathrm{A}} \mathrm{N}) \xleftarrow {\psi_ {\mathrm{M}} (\mathrm{N})} \operatorname{Tor} ^ {\mathrm{A}} (\mathrm{M}, \mathrm{N}) \xrightarrow {\overline {{\psi}} _ {\mathrm{N}} (\mathrm{M})} \mathrm{H} (\mathrm{M} \otimes_ {\mathrm{A}} \mathrm{L} (\mathrm{N}));
\]

根据X，第67页，命题4， \(1 \otimes p_{\mathbf{N}}\) 与 \(p_{\mathbf{M}} \otimes 1\) 是拟同构。因此：

命题5. --- k-同态

\[
\begin{array}{l} \psi_ {\mathrm{M}} (\mathrm{N}): \operatorname{Tor} ^ {\mathrm{A}} (\mathrm{M}, \mathrm{N}) \to \mathrm{H} (\mathrm{L} (\mathrm{M}) \otimes_ {\mathrm{A}} \mathrm{N}) \\ \overline {{\psi}} _ {\mathrm{N}} (\mathrm{M}): \operatorname{Tor} ^ {\mathrm{A}} (\mathrm{M}, \mathrm{N}) \to \mathrm{H} (\mathrm{M} \otimes_ {\mathrm{A}} \mathrm{L} (\mathrm{N})) \end{array}
\]

都是双射的。

推论。—— 若 M 或 N 是平坦的，则 \(\operatorname{Tor}_{i}^{\mathbf{A}}(\mathbf{M},\mathbf{N}) = 0\) 对 \(i \geqslant 0\) 成立。

假设 N（相应地 M）是平坦的；则 \(p_{\mathbf{M}} \otimes 1 : \mathbf{L}(\mathbf{M}) \otimes_{\mathbf{A}} \mathbf{N} \to \mathbf{M} \otimes_{\mathbf{A}} \mathbf{N}\) (相应地， \(1 \otimes p_{\mathbf{N}} : \mathbf{M} \otimes_{\mathbf{A}} \mathbf{L}(\mathbf{N}) \to \mathbf{M} \otimes_{\mathbf{A}} \mathbf{N}\) ）是一个拟同构（X，第 67 页，命题 4），因此 \(\mathrm{H}_{i}(\mathrm{L}(\mathrm{M}) \otimes_{\mathbf{A}} \mathrm{N})\) (相应地， \(\mathrm{H}_{i}(\mathrm{M} \otimes_{\mathbf{A}} \mathrm{L}(\mathrm{N}))\) ）对 i \textgreater{} 0 为零。

注记 3。—— 若 \(g: \mathbf{N} \to \mathbf{N}'\) 是左 A-模的同态，则

\[
\left(1 _ {\mathrm{L} (\mathbf {M})} \otimes g\right) \circ \left(1 _ {\mathrm{L} (\mathbf {M})} \otimes 1 _ {\mathrm{N}}\right) = \left(1 _ {\mathrm{L} (\mathbf {M})} \otimes 1 _ {\mathrm{N}}\right) \circ \left(1 _ {\mathrm{L} (\mathbf {M})} \otimes \mathrm{L} (g)\right),
\]

因此图表

\[
\begin{array}{c c c} \operatorname{Tor} ^ {\mathbf {A}} (\mathbf {M}, \mathbf {N}) & \xrightarrow {\psi_ {\mathbf {M}} (\mathbf {N})} & \mathrm{H} (\mathrm{L} (\mathbf {M}) \otimes_ {\mathbf {A}} \mathbf {N}) \\ \operatorname{Tor} ^ {\mathbf {A}} (1, g) \Big \downarrow & & \Big \downarrow \mathrm{H} (1 \otimes g) \\ \operatorname{Tor} ^ {\mathbf {A}} (\mathbf {M}, \mathbf {N} ^ {\prime}) & \xrightarrow {\psi_ {\mathbf {M}} (\mathbf {N} ^ {\prime})} & \mathrm{H} (\mathrm{L} (\mathbf {M}) \otimes_ {\mathbf {A}} \mathbf {N} ^ {\prime}) \end{array}
\]

是交换的。

同样地，若 \(f: \mathbf{M} \to \mathbf{M}'\) 是右 A-模的同态，我们有交换图：

\[
\begin{array}{c c} \operatorname{Tor} ^ {\mathbf {A}} (\mathrm{M}, \mathrm{N}) & \xrightarrow {\overline {{\psi}} _ {\mathrm{N}} (\mathrm{M})} \mathrm{H} (\mathrm{M} \otimes_ {\mathrm{A}} \mathrm{L} (\mathrm{N})) \\ \operatorname{Tor} ^ {\mathbf {A}} (f, 1) \Big \downarrow & \Big \downarrow \mathrm{H} (f \otimes 1) \\ \operatorname{Tor} ^ {\mathbf {A}} (\mathrm{M} ^ {\prime}, \mathrm{N}) & \xrightarrow {\overline {{\psi}} _ {\mathrm{N}} (\mathrm{M} ^ {\prime})} \mathrm{H} (\mathrm{M} ^ {\prime} \otimes_ {\mathrm{A}} \mathrm{L} (\mathrm{N})). \end{array}
\]

命题6. --- 映射 (f, g) ↦ Tor \(^{A}\) (f, g) :

\[
\operatorname{Hom} _ {\mathrm{A}} \left(\mathrm{M}, \mathrm{M} ^ {\prime}\right) \times \operatorname{Hom} _ {\mathrm{A}} \left(\mathrm{N}, \mathrm{N} ^ {\prime}\right)\rightarrow \operatorname{Hom} _ {k} \left(\operatorname{Tor} ^ {\mathrm{A}} (\mathrm{M}, \mathrm{N}), \operatorname{Tor} ^ {\mathrm{A}} \left(\mathrm{M} ^ {\prime}, \mathrm{N} ^ {\prime}\right)\right)
\]

是 k-双线性的。

设 \(f \in \operatorname{Hom}_{\mathbf{A}}(\mathbf{M}, \mathbf{M}')\) , \(g_1, g_2 \in \operatorname{Hom}_{\mathbf{A}}(\mathbf{N}, \mathbf{N}')\) , \(\lambda_1, \lambda_2 \in k\) 。则态射

\[
\lambda_ {1} (\mathrm{L} (f) \otimes g _ {1}) + \lambda_ {2} (\mathrm{L} (f) \otimes g _ {2}) \quad \text { and } \quad \mathrm{L} (f) \otimes (\lambda_ {1} g _ {1} + \lambda_ {2} g _ {2})
\]

这两个从 L(M) ⊗\_A N 到 L(M) ⊗\_A N′ 的态射相一致；由命题5和注记3，我们因此有

\[
\operatorname{Tor} ^ {\mathbf {A}} \left(f, \lambda_ {1} g _ {1} + \lambda_ {2} g _ {2}\right) = \lambda_ {1} \operatorname{Tor} ^ {\mathbf {A}} \left(f, g _ {1}\right) + \lambda_ {2} \operatorname{Tor} ^ {\mathbf {A}} \left(f, g _ {2}\right).\tag{10}
\]

我们对映射 \(f\mapsto\mathrm{Tor}^{\mathbf{A}}(f,g)\) 作类似推理。

推论. --- 设 \(\lambda\in k\) 。若 \(\lambda\) 湮灭 M 或 N，则它湮灭 Tor \(^{A}\) (M, N)。

事实上， \(\lambda .1_{\mathrm{Tor}(\mathbf{M},\mathbf{N})} = \mathrm{Tor}(\lambda .1_{\mathbf{M}},1_{\mathbf{N}}) = \mathrm{Tor}(1_{\mathbf{M}},\lambda .1_{\mathbf{N}})\) .

命题7. --- 设I和J为两个集合， \((\mathbf{M}_{\alpha})_{\alpha \in I}\) 一族右 A-模， \((\mathbf{N}_{\beta})_{\beta \in J}\) 一个左A-模族。同态

\[
\bigoplus_ {\alpha \in I, \beta \in J} \operatorname{Tor} ^ {A} \left(M _ {\alpha}, N _ {\beta}\right)\rightarrow \operatorname{Tor} ^ {A} \left(\bigoplus_ {\alpha \in I} M _ {\alpha}, \bigoplus_ {\beta \in J} N _ {\beta}\right)
\]

由典范同态 \(M_{\alpha} \rightarrow \oplus M_{\alpha}\) 与 \(N_{\beta} \rightarrow \oplus N_{\beta}\) 导出的，是双射的。只需证明对每个右模M（相应地，每个左模N），典范同态

\[
\bigoplus_ {\beta \in J} \operatorname{Tor} ^ {A} (M, N _ {\beta}) \rightarrow \operatorname{Tor} ^ {A} (M, \bigoplus_ {\beta \in J} N _ {\beta})
\]

\((\text{resp. } \bigoplus_{\alpha \in I} \text{Tor}^{\mathbf{A}}(\mathbf{M}_{\alpha}, \mathbf{N}) \to \text{Tor}^{\mathbf{A}} (\bigoplus_{\alpha \in I} \mathbf{M}_{\alpha}, \mathbf{N}))\) 是双射的。这由前述内容、X，第28页，命题1以及以下典范同构得出：

\[
\begin{array}{l} \underset {\beta} {\oplus} (\mathrm{L(M)} \otimes_ {\mathrm{A}} \mathrm {N}_ {\beta}) \to \mathrm{L(M)} \otimes_ {\mathrm{A}} (\underset {\beta} {\oplus} \mathrm {N}_ {\beta}), \\ \underset {\alpha} {\oplus} (\mathrm {M}_ {\alpha} \otimes_ {\mathrm{A}} \mathrm{L(N)}) \to (\underset {\alpha} {\oplus} \mathrm {M}_ {\alpha}) \otimes_ {\mathrm{A}} \mathrm{L(N)}. \end{array}
\]

类似的论证给出：

命题8.——设I（相应地，J）是一个右有向预序集， \(((M_{\alpha}), (u_{\alpha' \alpha}))\) (相应地， \(((N_{\beta}), (v_{\beta' \beta}))\) ) 是关于 I（相应地，J）的右（相应地，左）A-模的归纳系。分次 k-模的同态

\[
\lim _ {(\alpha , \beta) \in I \times J} \operatorname{Tor} ^ {A} \left(M _ {\alpha}, N _ {\beta}\right)\rightarrow \operatorname{Tor} ^ {A} \left(\lim _ {\alpha \in I} M _ {\alpha}, \lim _ {\beta \in J} N _ {\beta}\right),
\]

由典范的 A-同态 \(\mathbf{M}_{\alpha} \to \varinjlim \mathbf{M}_{\alpha}\) 与 \(\mathbf{N}_{\beta} \to \varinjlim \mathbf{N}_{\beta}\) 导出的，是双射。

特别地，取 \(J = I\) 并注意到 \((\alpha, \alpha), \alpha \in I\) 构成 \(I \times I\) 的共尾子集，我们得到：

推论. --- 设 I 是一个右有向预序集， \((\mathbf{M}_{i}, u_{ji})\) (相应地， \((\mathbf{N}_{i}, v_{ji})\) ) 是关于 I 的右（相应地，左）A-模的归纳系。分次 k-模的同态

\[
\varinjlim_ {i \in I} \operatorname{Tor} ^ {A} (M _ {i}, N _ {i}) \rightarrow \operatorname{Tor} ^ {A} (\varinjlim_ {i \in I} M _ {i}, \varinjlim_ {i \in I} N _ {i}),
\]

由典范的 A-同态 \(\mathbf{M}_i\to \varinjlim \mathbf{M}_i\) 与 \(\mathbf{N}_j\to \varinjlim \mathbf{N}_j\) 导出的，是双射。

设 M 为右 A-模，N 为左 A-模，A° 为 A 的反环，M° 为 M 所对应的左 A°-模，N° 为 N 所对应的右 A°-模。我们有 L(M°) = L(M)° 且 L(N°) = L(N)°，从而得到一个交换同构（X，第 63 页，命题 2）

\[
\sigma (\mathrm{L} (\mathrm{M}), \mathrm{L} (\mathrm{N})): \mathrm{L} (\mathrm{M}) \otimes_ {\mathrm{A}} \mathrm{L} (\mathrm{N}) \rightarrow \mathrm{L} (\mathrm{N} ^ {\circ}) \otimes_ {\mathrm{A} ^ {0}} \mathrm{L} (\mathrm{M} ^ {\circ}).
\]

过渡到同调， \(\sigma(\mathrm{L}(\mathbf{M}),\mathrm{L}(\mathbf{N}))\) 诱导出一个次数为 0 的分次同构 \(\sigma_{\mathbf{M},\mathbf{N}}:\operatorname{Tor}^{\mathbf{A}}(\mathbf{M},\mathbf{N})\to\operatorname{Tor}^{\mathbf{A}^{\circ}}(\mathbf{N}^{\circ},\mathbf{M}^{\circ})\) 称为挠积的交换同构。

注意 \(\sigma_{\mathbf{N}^{\circ},\mathbf{M}^{\circ}}\circ\sigma_{\mathbf{M},\mathbf{N}}=\mathrm{Id}_{\mathrm{Tor}\left(\mathbf{M},\mathbf{N}\right)}\) 并且 \(\sigma_{\mathbf{M},\mathbf{N}}\) 在零次项上诱导出张量积的交换同态。另一方面，若 \(f:\mathbf{M}\to\mathbf{M}'\) 与 \(g:\mathbf{N}\to\mathbf{N}'\) 是A-模的同态，我们有

\[
\operatorname{Tor} ^ {\mathbf {A} ^ {\circ}} (g, f) \circ \sigma_ {\mathbf {M}, \mathbf {N}} = \sigma_ {\mathbf {M} ^ {\prime}, \mathbf {N} ^ {\prime}} \circ \operatorname{Tor} ^ {\mathbf {A}} (f, g).
\]

\subsection{连接同态与正合序列}

设M是一个右A-模。回忆对于每个左A-模N，我们在前一小节（X，第69页，命题5）中定义了一个同构

设

\[
\psi_ {\mathbf {M}} (\mathbf {N}): \operatorname{Tor} ^ {\mathbf {A}} (\mathbf {M}, \mathbf {N}) \rightarrow \mathrm{H} (\mathrm{L} (\mathbf {M}) \otimes_ {\mathbf {A}} \mathbf {N}).\tag{ε}
\]

\[
0 \rightarrow \mathbf {N} ^ {\prime} \xrightarrow {u} \mathbf {N} \xrightarrow {v} \mathbf {N} ^ {\prime \prime} \rightarrow 0
\]

一个左A-模的正合序列；则 \(k\) -模复形的序列

\[
\left(^ {\mathrm{M}} \mathcal {E}\right) \quad 0 \longrightarrow \mathrm{L} (\mathrm{M}) \otimes_ {\mathrm{A}} \mathrm{N} ^ {\prime} \xrightarrow {1 \otimes u} \mathrm{L} (\mathrm{M}) \otimes_ {\mathrm{A}} \mathrm{N} \xrightarrow {1 \otimes v} \mathrm{L} (\mathrm{M}) \otimes_ {\mathrm{A}} \mathrm{N} ^ {\prime \prime} \longrightarrow 0
\]

这一序列是正合的（X，第66页，引理1）；令

\[
\hat {c} \left(^ {\mathrm{M}} \mathcal {E}\right): \mathrm{H} (\mathrm{L} (\mathrm{M}) \otimes_ {\mathrm{A}} \mathrm{N} ^ {\prime \prime}) \rightarrow \mathrm{H} (\mathrm{L} (\mathrm{M}) \otimes_ {\mathrm{A}} \mathrm{N} ^ {\prime})
\]

为相应的连接同态（X，第 29 页）。

定义 2. --- 相对于模 M 和正合序列 \(\mathcal{E}\) 的挠积的连接同态指复合同态

\[
\partial (\mathrm{M}, \mathcal {E}) = \psi_ {\mathrm{M}} \left(\mathrm{N} ^ {\prime}\right) ^ {- 1} \circ \partial \left(^ {\mathrm{M}} \mathcal {E}\right) \circ \psi_ {\mathrm{M}} \left(\mathrm{N} ^ {\prime \prime}\right): \operatorname{Tor} ^ {\mathrm{A}} \left(\mathrm{M}, \mathrm{N} ^ {\prime \prime}\right)\rightarrow \operatorname{Tor} ^ {\mathrm{A}} \left(\mathrm{M}, \mathrm{N} ^ {\prime}\right).
\]

这是一个分次 \(k\) -同态，次数为 \((-1)\) ，其齐次分量记为 \(\partial_n(\mathbf{M}, \mathcal{E}) : \mathrm{Tor}_n^{\mathbf{A}}(\mathbf{M}, \mathbf{N}'') \to \mathrm{Tor}_{n-1}^{\mathbf{A}}(\mathbf{M}, \mathbf{N}')\) .

定理 1. --- k-模同态的向左无界序列

\[
\begin{array}{l} \dots \longrightarrow \operatorname{Tor} _ {n} ^ {\mathbf {A}} (\mathbf {M}, \mathbf {N} ^ {\prime}) \xrightarrow {\operatorname{Tor} _ {n} ^ {\mathbf {A}} (1 , u)} \operatorname{Tor} _ {n} ^ {\mathbf {A}} (\mathbf {M}, \mathbf {N}) \xrightarrow {\operatorname{Tor} _ {n} ^ {\mathbf {A}} (1 , v)} \operatorname{Tor} _ {n} ^ {\mathbf {A}} (\mathbf {M}, \mathbf {N} ^ {\prime \prime}) \\ \xrightarrow {\hat {c} _ {n} (\mathbf {M} , \mathcal {E})} \operatorname{Tor} _ {n - 1} ^ {\mathbf {A}} (\mathbf {M}, \mathbf {N} ^ {\prime}) \xrightarrow {\operatorname{Tor} _ {n - 1} ^ {\mathbf {A}} (1 , u)} \dots \xrightarrow {\operatorname{Tor} _ {1} ^ {\mathbf {A}} (1 , v)} \operatorname{Tor} _ {1} ^ {\mathbf {A}} (\mathbf {M}, \mathbf {N} ^ {\prime \prime}) \\ \xrightarrow {\hat {c} _ {1} (\mathbf {M} , \mathcal {E})} \mathbf {M} \otimes_ {\mathbf {A}} \mathbf {N} ^ {\prime} \xrightarrow {1 \otimes u} \mathbf {M} \otimes_ {\mathbf {A}} \mathbf {N} \xrightarrow {1 \otimes v} \mathbf {M} \otimes_ {\mathbf {A}} \mathbf {N} ^ {\prime \prime} \longrightarrow 0 \end{array}
\]

是正合的。

事实上，考虑图表

\[
\begin{array}{c} \operatorname{Tor} (\mathbf {M}, \mathbf {N} ^ {\prime}) \xrightarrow {\operatorname{Tor} (1 , u)} \operatorname{Tor} (\mathbf {M}, \mathbf {N}) \xrightarrow {\operatorname{Tor} (1 , v)} \operatorname{Tor} (\mathbf {M}, \mathbf {N} ^ {\prime \prime}) \xrightarrow {\partial (\mathbf {M} , \mathcal {E})} \operatorname{Tor} (\mathbf {M}, \mathbf {N} ^ {\prime}) \xrightarrow {\operatorname{Tor} (1 , u)} \operatorname{Tor} (\mathbf {M}, \mathbf {N}) \\ \psi_ {\mathbf {M} (\mathbf {N} ^ {\prime})} \Big | _ {\downarrow} \quad \psi_ {\mathbf {M} (\mathbf {N})} \Big | _ {\downarrow} \quad \psi_ {\mathbf {M} (\mathbf {N} ^ {\prime \prime})} \Big | _ {\downarrow} \quad \psi_ {\mathbf {M} (\mathbf {N} ^ {\prime})} \Big | _ {\downarrow} \quad \psi_ {\mathbf {M} (\mathbf {N} ^ {\prime})} \Big | _ {\downarrow} \\ H (L (\mathbf {M}) \otimes N ^ {\prime}) \xrightarrow {H (1 \otimes u)} H (L (\mathbf {M}) \otimes N) \xrightarrow {H (1 \otimes v)} H (L (\mathbf {M} \otimes N ^ {\prime \prime}) \xrightarrow {\partial^ {(M)} \mathcal {E}} H (L (\mathbf {M}) \otimes N ^ {\prime}) \xrightarrow {} H (L (\mathbf {M}) \otimes N). \end{array}
\]

由（X，第 69 页，注记 3）和定义 2 知它是交换的。另一方面，底行是正合的（X，第 30 页，定理 1），并且各个 \(\psi_{\mathrm{M}}\) 是双射的（X，第 69 页，命题 5）。

推论 1. --- 如果 Tor \(_{1}^{A}\) (M, N'') = 0，则序列

\[
0 \longrightarrow \mathrm{M} \otimes_ {\mathrm{A}} \mathrm{N} ^ {\prime} \xrightarrow {1 \otimes u} \mathrm{M} \otimes_ {\mathrm{A}} \mathrm{N} \xrightarrow {1 \otimes v} \mathrm{M} \otimes_ {\mathrm{A}} \mathrm{N} ^ {\prime \prime} \longrightarrow 0
\]

是正合的。

推论 2. --- 设 \(0 \to C' \xrightarrow{u} C \xrightarrow{v} C'' \to 0\) 是左 A-模复形的正合序列，且 E 是右 A-模复形。如果 \(C''\) 或 E 是平坦的，则序列

\[
0 \longrightarrow \mathrm{E} \otimes_ {\mathrm{A}} \mathrm{C} ^ {\prime} \xrightarrow {1 \otimes u} \mathrm{E} \otimes_ {\mathrm{A}} \mathrm{C} \xrightarrow {1 \otimes v} \mathrm{E} \otimes_ {\mathrm{A}} \mathrm{C} ^ {\prime \prime} \longrightarrow 0
\]

是正合的。

事实上， \(\operatorname{Tor}_1^{\mathbf{A}}(\mathbf{E},\mathbf{C}^{\prime \prime}) = 0\) 根据 X，第 69 页，命题 5 的推论。

例. --- 设 a 是 A 的一个理想。左 A-模的正合序列

\[
0 \rightarrow \mathfrak {a} \rightarrow \mathrm{A} _ {\mathrm{s}} \rightarrow \mathrm{A} / \mathfrak {a} \rightarrow 0
\]

给出挠积的正合序列，其中项 \(\operatorname{Tor}_{i}^{\mathbf{A}}(\mathbf{M},\mathbf{A})\) 对于 i \textgreater{} 0 为零。由此我们推出同构

\[
\operatorname{Tor} _ {i + 1} ^ {\mathbf {A}} (\mathbf {M}, \mathbf {A} / \mathfrak {a}) \rightarrow \operatorname{Tor} _ {i} ^ {\mathbf {A}} (\mathbf {M}, \mathfrak {a}), \quad i > 0
\]

以及一个正合序列

\[
0 \rightarrow \operatorname{Tor} _ {1} ^ {\mathbf {A}} (\mathrm{M}, \mathrm{A} / \mathfrak {a}) \rightarrow \mathrm{M} \otimes_ {\mathbf {A}} \mathfrak {a} \rightarrow \mathrm{M} \otimes_ {\mathbf {A}} \mathrm{A} \rightarrow \mathrm{M} \otimes \mathrm{A} / \mathfrak {a} \rightarrow 0:
\]

由此可知 \(\operatorname{Tor}_{1}^{\mathbf{A}}(\mathbf{M},\mathbf{A}/\mathfrak{a})\) 与典范同态 \(\mathbf{M}\otimes_{\mathbf{A}}\mathfrak{a}\to\mathbf{M}\) 的核等同。

例如，取M为如下形式的模： \(A_{d}/b\) ，其中 b 是 A 的一个右理想，则得到 \(\operatorname{Tor}_{1}^{A}(A/b, A/a)\) 到 \((a \cap b)/ba\) 上的一个同构。

命题 9. --- 设 \(f: \mathbf{M} \to \mathbf{M}_{1}\) 是右A-模的同态，且

(ε)

\[
\begin{array}{c} 0 \longrightarrow \mathrm{N} ^ {\prime} \stackrel {{u}} {{\longrightarrow}} \mathrm{N} \stackrel {{v}} {{\longrightarrow}} \mathrm{N} ^ {\prime \prime} \longrightarrow 0 \\ g ^ {\prime} \Biggl \downarrow \qquad \qquad \qquad \qquad \qquad \qquad \qquad \qquad \qquad \qquad \qquad \qquad \qquad \qquad \qquad \qquad \qquad \qquad \qquad \qquad \qquad \qquad \qquad \qquad \qquad \qquad \qquad \qquad \qquad \qquad \qquad \qquad \qquad \qquad \\ 0 \longrightarrow \mathrm{N} _ {1} ^ {\prime} \stackrel {{u _ {1}}} {{\longrightarrow}} \mathrm{N} _ {1} \stackrel {{v _ {1}}} {{\longrightarrow}} \mathrm{N} _ {1} ^ {\prime \prime} \longrightarrow 0 \end{array}\tag{$(\mathcal{E}_1)$}
\]

一个具有正合行的左 A-模同态的交换图。k-模的图

\[
\begin{array}{c} \operatorname{Tor} ^ {\mathbf {A}} (\mathbf {M}, \mathbf {N} ^ {\prime \prime}) \xrightarrow {\hat {c} (\mathbf {M} , \mathcal {E})} \operatorname{Tor} ^ {\mathbf {A}} (\mathbf {M}, \mathbf {N} ^ {\prime}) \\ \operatorname{Tor} ^ {\mathbf {A}} (f, g ^ {\prime \prime}) \Big \downarrow \\ \operatorname{Tor} ^ {\mathbf {A}} (\mathbf {M} _ {1}, \mathbf {N} _ {1} ^ {\prime \prime}) \xrightarrow {\hat {c} (\mathbf {M} _ {1} , \mathcal {E} _ {1})} \operatorname{Tor} ^ {\mathbf {A}} (\mathbf {M} _ {1}, \mathbf {N} _ {1} ^ {\prime}) \end{array}
\]

是交换的。

这由将 X，第 31 页，命题 2 应用于该交换图得出

\[
\begin{array}{l} 0 \longrightarrow \mathrm{L(M)} \otimes_ {\mathrm{A}} \mathrm{N} ^ {\prime} \xrightarrow {1 \otimes u} \mathrm{L(M)} \otimes_ {\mathrm{A}} \mathrm{N} \xrightarrow {1 \otimes v} \mathrm{L(M)} \otimes_ {\mathrm{A}} \mathrm{N} ^ {\prime \prime} \longrightarrow 0 \\ \quad \mathrm{L(f)} \otimes g ^ {\prime} \Biggl \downarrow \quad \mathrm{L(f)} \otimes g \Biggl \downarrow \quad \mathrm{L(f)} \otimes g ^ {\prime \prime} \Biggl \downarrow \\ 0 \longrightarrow \mathrm{L(M} _ {1}) \otimes_ {\mathrm{A}} \mathrm{N} _ {1} ^ {\prime} \xrightarrow {1 \otimes u _ {1}} \mathrm{L(M} _ {1}) \otimes_ {\mathrm{A}} \mathrm{N} _ {1} \xrightarrow {1 \otimes v _ {1}} \mathrm{L(M} _ {1}) \otimes_ {\mathrm{A}} \mathrm{N} _ {1} ^ {\prime \prime} \longrightarrow 0. \end{array}
\]

类似地，若 N 是一个左 A-模且

\[
0 \to \mathbf {M} ^ {\prime} \xrightarrow {r} \mathbf {M} \xrightarrow {s} \mathbf {M} ^ {\prime \prime} \to 0\tag{F}
\]

一个右 A-模的正合序列，定义连接同态

\[
\partial (\mathcal {F}, \mathrm{N}): \operatorname{Tor} ^ {\mathrm{A}} \left(\mathrm{M} ^ {\prime \prime}, \mathrm{N}\right)\rightarrow \operatorname{Tor} ^ {\mathrm{A}} \left(\mathrm{M} ^ {\prime}, \mathrm{N}\right)
\]

\[
\partial_ {n} (\mathcal {F}, \mathrm{N}): \operatorname{Tor} _ {n} ^ {\mathrm{A}} \left(\mathrm{M} ^ {\prime \prime}, \mathrm{N}\right)\rightarrow \operatorname{Tor} _ {n - 1} ^ {\mathrm{A}} \left(\mathrm{M} ^ {\prime}, \mathrm{N}\right)
\]

由 \(\partial(\mathcal{F},\mathrm{N})=\overline{\psi}_{\mathrm{N}}(\mathrm{M}^{\prime})^{-1}\circ\partial(\mathcal{F}^{\mathrm{N}})\circ\overline{\psi}_{\mathrm{N}}(\mathrm{M}^{\prime\prime})\) 定义，其中 \(\partial(\mathcal{F}^{\mathrm{N}})\) 是正合序列的连接同态

\[
(\mathcal {F} ^ {N})
\]

\[
0 \rightarrow \mathrm{M} ^ {\prime} \otimes_ {\mathrm{A}} \mathrm{L} (\mathrm{N}) \rightarrow \mathrm{M} \otimes_ {\mathrm{A}} \mathrm{L} (\mathrm{N}) \rightarrow \mathrm{M} ^ {\prime \prime} \otimes_ {\mathrm{A}} \mathrm{L} (\mathrm{N}) \rightarrow 0
\]

由 \(\mathcal{F}\) 导出的，并且我们有：

定理 1 bis. --- k-模同态的向左无界序列

\[
\begin{array}{l}\rightarrow \operatorname{Tor} _ {n} ^ {\mathsf {A}} (\mathbf {M} ^ {\prime}, \mathbf {N}) \xrightarrow {\operatorname{Tor} _ {n} ^ {\mathsf {A}} (r , 1)} \operatorname{Tor} _ {n} ^ {\mathsf {A}} (\mathbf {M}, \mathbf {N}) \xrightarrow {\operatorname{Tor} _ {n} ^ {\mathsf {A}} (s , 1)} \operatorname{Tor} _ {n} ^ {\mathsf {A}} (\mathbf {M} ^ {\prime \prime}, \mathbf {N}) \xrightarrow {\partial_ {n} (\mathcal {F} , \mathbf {N})} \operatorname{Tor} _ {n - 1} ^ {\mathsf {A}} (\mathbf {M} ^ {\prime}, \mathbf {N})\\\dots \rightarrow \operatorname{Tor} _ {1} ^ {\mathsf {A}} (\mathbf {M} ^ {\prime \prime}, \mathbf {N}) \xrightarrow {\partial_ {1} (\mathcal {F} , \mathbf {N})} \mathbf {M} ^ {\prime} \otimes_ {\mathbf {A}} \mathbf {N} \xrightarrow {r \otimes 1} \mathbf {M} \otimes_ {\mathbf {A}} \mathbf {N} \xrightarrow {s \otimes 1} \mathbf {M} ^ {\prime \prime} \otimes_ {\mathbf {A}} \mathbf {N} \longrightarrow 0\end{array}
\]

是正合的。

我们留给读者来陈述并证明与定理1的推论及命题9相类似的性质。此外：

命题10。——设( \(F^{\circ}\) ) 表示左 A-模的正合序列

\[
0 \rightarrow \mathrm{M} ^ {\prime \circ} \xrightarrow {r} \mathrm{M} ^ {\circ} \xrightarrow {s} \mathrm{M} ^ {\prime \prime \circ} \rightarrow 0.
\]

图表

\[
\begin{array}{c} \operatorname{Tor} ^ {\mathbf {A}} (\mathbf {M} ^ {\prime \prime}, \mathbf {N}) \xrightarrow {\partial (\mathcal {F} , \mathbf {N})} \operatorname{Tor} ^ {\mathbf {A}} (\mathbf {M} ^ {\prime}, \mathbf {N}) \\ \sigma_ {\mathbf {M} ^ {\prime}, \mathbf {N}} \Biggl \downarrow \qquad \qquad \qquad \qquad \qquad \qquad \qquad \qquad \qquad \sigma_ {\mathbf {M} ^ {\prime}, \mathbf {N}} \Biggl \downarrow \\ \operatorname{Tor} ^ {\mathbf {A} ^ {\circ}} (\mathbf {N} ^ {\circ}, \mathbf {M} ^ {\prime \prime^ {\circ}}) \xrightarrow {\partial (\mathbf {N} ^ {\circ} , \mathcal {F} ^ {\circ})} \operatorname{Tor} ^ {\mathbf {A} ^ {\circ}} (\mathbf {N} ^ {\circ}, \mathbf {M} ^ {\prime \circ}) \end{array}
\]

是交换的。

事实上，这由将 X，第 31 页，命题 2 应用于交换图表得出

\[
\begin{array}{l} 0 \longrightarrow \mathrm{M} ^ {\prime} \otimes_ {\mathrm{A}} \mathrm{L} (\mathrm{N}) \xrightarrow {r \otimes 1} \mathrm{M} \otimes_ {\mathrm{A}} \mathrm{L} (\mathrm{N}) \xrightarrow {s \otimes 1} \mathrm{M} ^ {\prime \prime} \otimes_ {\mathrm{A}} \mathrm{L} (\mathrm{N}) \longrightarrow 0 \\ \sigma (\mathrm{M} ^ {\prime}, \mathrm{L} (\mathrm{N})) \Biggl \downarrow \quad \sigma (\mathrm{M}, \mathrm{L} (\mathrm{N})) \Biggl \downarrow \quad \sigma (\mathrm{M} ^ {\prime \prime}, \mathrm{L} (\mathrm{N})) \Biggl \downarrow \\ 0 \longrightarrow \mathrm{L} (\mathrm{N} ^ {\circ}) \otimes_ {\mathrm{A} ^ {\circ}} \mathrm{M} ^ {\prime \circ} \xrightarrow {1 \otimes r} \mathrm{L} (\mathrm{N} ^ {\circ}) \otimes_ {\mathrm{A} ^ {\circ}} \mathrm{M} ^ {\circ} \xrightarrow {1 \otimes s} \mathrm{L} (\mathrm{N} ^ {\circ}) \otimes_ {\mathrm{A} ^ {\circ}} \mathrm{M} ^ {\prime \prime \circ} \longrightarrow 0. \end{array}
\]

我们将在后面看到其他的交换关系（参见 X，第 131 页，推论 1）。

\subsection{平坦模与挠积}

定理 2. —— 设 E 是一个右 A-模。下列条件是等价的：

\begin{enumerate}
\def\labelenumi{(\roman{enumi})}
\item
  E 是平坦的；
\item
  对于每个左 A-模 F 以及每个整数 n \textgreater{} 0，有
\end{enumerate}

\[
\operatorname{Tor} _ {n} ^ {\mathbf {A}} (\mathbf {E}, \mathbf {F}) = 0;
\]

\begin{enumerate}
\def\labelenumi{(\roman{enumi})}
\setcounter{enumi}{2}
\tightlist
\item
  对于每个有限呈现的单生成左 A-模 F，有
\end{enumerate}

\[
\operatorname{Tor} _ {1} ^ {\mathbf {A}} (\mathbf {E}, \mathbf {F}) = 0;
\]

\begin{enumerate}
\def\labelenumi{(\roman{enumi})}
\setcounter{enumi}{3}
\item
  对于 A 的每个有限生成左理想 a，典范映射 E ⊗\_A a → E 是单射；
\item
  对于如下形式的右 A-模的每个正合序列
\end{enumerate}

\[
0 \rightarrow \mathbf {G} \xrightarrow {v} \mathbf {H} \xrightarrow {w} \mathbf {E} \rightarrow 0,
\]

以及每个左 A-模 F，序列

\[
0 \longrightarrow \mathrm{G} \otimes_ {\mathrm{A}} \mathrm{F} \xrightarrow {v \otimes 1} \mathrm{H} \otimes_ {\mathrm{A}} \mathrm{F} \xrightarrow {w \otimes 1} \mathrm{E} \otimes_ {\mathrm{A}} \mathrm{F} \longrightarrow 0
\]

是正合的。

\begin{enumerate}
\def\labelenumi{(\roman{enumi})}
\item
  \(\Rightarrow\) （ii）：这是 X，第 69 页，命题 5 的推论。
\item
  \(\Rightarrow\) （iii）：这是平凡的。
\item
  \(\Leftrightarrow\) （iv）：每个有限呈现的单生成左 A-模都同构于一个商模 A/a，其中 a 是一个有限生成的左理想，因此由 X，第 72 页，例可知（iii）等价于（iv）。
\item
  \(\Rightarrow\) （i）：由 X，第 8 页，命题 3，X，第 72 页，定理 1，只要 \(\operatorname{Tor}_{1}^{\mathbf{A}}(\mathbf{E},\mathbf{F}) = 0\) 对每个左 A-模 F 成立，E 就是平坦的。若（iii）成立，则当 F 是单生成且有限呈现时，上述条件成立。由 X，第 11 页，命题 7，每个 A-模（相应地，每个单生成 A-模）都是有限呈现模（相应地，有限呈现的单生成模）的滤过归纳极限；因此，由 X，第 70 页，命题 8，可以看出只需证明：若 \(\operatorname{Tor}_{1}^{\mathbf{A}}(\mathbf{E},\mathbf{F}) = 0\) 当 F 是单生成时成立，则当 F 是有限生成时也成立。于是我们对 F 的生成族 \((f_{1},\ldots,f_{n})\) 的基数进行归纳；正合序列
\end{enumerate}

\[
0 \rightarrow \mathrm{A} f _ {1} \rightarrow \mathrm{F} \rightarrow \mathrm{F} / \mathrm{A} f _ {1} \rightarrow 0
\]

给出正合序列

\[
\operatorname{Tor} _ {1} ^ {\mathbf {A}} (\mathrm{E}, \mathrm{A} f _ {1}) \rightarrow \operatorname{Tor} _ {1} ^ {\mathbf {A}} (\mathrm{E}, \mathrm{F}) \rightarrow \operatorname{Tor} _ {1} ^ {\mathbf {A}} (\mathrm{E}, \mathrm{F} / \mathrm{A} f _ {1}),
\]

从而 \(\operatorname{Tor}_{1}^{\mathbf{A}}(\mathbf{E},\mathbf{F})=0\)，因为 \(\operatorname{Tor}_{1}^{\mathbf{A}}(\mathbf{E},\mathbf{A}f_{1})=0\) 且 \(\operatorname{Tor}_{1}^{\mathbf{A}}(\mathbf{E},\mathbf{F}/\mathbf{A}f_{1})=0\)，由归纳假设。

\begin{enumerate}
\def\labelenumi{(\roman{enumi})}
\item
  \(\Rightarrow\) (v)：这是定理1的推论2（X，第72页）。
\item
  \(\Rightarrow\) (iii)：正合序列（X，第50页）
\end{enumerate}

\[
0 \longrightarrow Z _ {0} (\mathrm{E}) \stackrel {i _ {\mathrm{E}}} {\longrightarrow} L _ {0} (\mathrm{E}) \stackrel {p _ {\mathrm{E}}} {\longrightarrow} \mathrm{E} \longrightarrow 0
\]

对每个左A-模F，产生一个正合序列

\[
0 \longrightarrow \operatorname{Tor} _ {1} ^ {\mathrm{A}} (\mathrm{E}, \mathrm{F}) \longrightarrow Z _ {0} (\mathrm{E}) \otimes_ {\mathrm{A}} \mathrm{F} \xrightarrow {i _ {\mathrm{E}} \otimes 1} \mathrm{L} _ {0} (\mathrm{E}) \otimes_ {\mathrm{A}} \mathrm{F} \xrightarrow {p _ {\mathrm{E}} \otimes 1} \mathrm{E} \otimes_ {\mathrm{A}} \mathrm{F} \longrightarrow 0.
\]

若满足(v)，则有 \(\operatorname{Tor}_1^{\mathbf{A}}(\mathbf{E},\mathbf{F}) = 0\)，由此得（iii）。

推论 1. --- 设 \(0 \to E' \to E \to E'' \to 0\) 是右A-模的一个正合序列。假设 \(E''\) 是平坦的。那么E平坦当且仅当 \(E'\) 是平坦的。

设 F 为左 A-模。由于 \(\operatorname{Tor}_i^{\mathbf{A}}(\mathbf{E}^{\prime \prime},\mathbf{F}) = 0\)（\(i = 1,2\)；定理2，(i) \(\Rightarrow\) (ii)），存在一个正合序列

\[
0 \to \mathrm{Tor} _ {1} ^ {\mathbf {A}} (\mathbf {E} ^ {\prime}, \mathbf {F}) \to \mathrm{Tor} _ {1} ^ {\mathbf {A}} (\mathbf {E}, \mathbf {F}) \to 0
\]

由此得出断言（定理2，（i）⇔（iii））。

推论 2. --- 设 \(0 \to E_n \to E_{n-1} \to \cdots \to E_1 \to 0\) 是右A-模的一个正合序列。如果 \(E_i\) 对 \(i = 1, \ldots, n-1\) 是平坦的，则 \(E_n\) 是平坦的。
